% concatenated from total.tex
\documentclass{article}
\usepackage{amsmath, amssymb, amsthm, graphicx, hyperref}
\usepackage{enumitem}
\usepackage{listings}
\usepackage{xcolor}
\usepackage{epstopdf}   % автоматическая конвертация .ps -> .pdf

\lstset{
  language=Python,
  basicstyle=\ttfamily\scriptsize,
  keywordstyle=\color{blue!70!black}\bfseries,
  commentstyle=\color{green!45!black}\itshape,
  stringstyle=\color{orange!70!black},
  numbers=left,
  numberstyle=\tiny\color{gray},
  numbersep=5pt,
  stepnumber=5,
  breaklines=true,
  breakatwhitespace=false,
  postbreak=\mbox{\textcolor{gray}{$\hookrightarrow$}\space},
  frame=single,
  columns=fullflexible,
  showstringspaces=false,
  tabsize=4,
  captionpos=b,
}

\newtheorem{theorem}{Theorem}[section]
\newtheorem{proposition}[theorem]{Proposition}
\newtheorem{corollary}[theorem]{Corollary}
\newtheorem{lemma}[theorem]{Lemma}
\newtheorem{example}[theorem]{Example}
\theoremstyle{remark}
\newtheorem{remark}[theorem]{Remark}
\newtheorem{definition}[theorem]{Definition}

\title{Operator Michelson Contrast: Logarithmic Parametrisation, Subadditivity, and a Sharp Noncommutativity Threshold}
\author{Irina Nikolaeva\thanks{Lean Technologies Laboratory, Kazan, Tatarstan, 420000, Krasnokokshayskaia 189a-202} \and Andrei Novikov\thanks{Novikov Laboratories, Kazan, Tatarstan, 420124, Absaliamova 14-124}}
\date{}

\providecommand{\QLC}{\operatorname{OLC}}
\providecommand{\QMC}{\operatorname{OMC}}
\providecommand{\MC}{\operatorname{MC}}
\providecommand{\Emax}{E_{\max}}
\providecommand{\Emin}{E_{\min}}

\begin{document}
\maketitle

\begin{abstract}
Building on the operator Michelson contrast introduced in our previous work for positive operators, we investigate its natural logarithmic parametrisation. For a positive invertible operator $A$, the operator Michelson contrast is $\Delta(A) = (\kappa(A)-1)/(\kappa(A)+1)$, where $\kappa(A) = \|A\|\|A^{-1}\|$ is the condition number. Its logarithmic lift is $\operatorname{arctanh}(\Delta(A)) = \tfrac12\ln\kappa(A)$. We call this quantity the Operator Logarithmic Contrast (OLC), and extend it to all invertible operators via the polar decomposition. We prove that QLC is subadditive under multiplication: $\QLC(AB)\le\QLC(A)+\QLC(B)$. For positive invertible operators, we establish an exact geometric identification: $\QLC(A)$ equals the projective Thompson distance from the identity ray to the ray generated by $A$. We prove a sharp dimensional threshold for equality in the subadditivity law: for $2\times2$ and $3\times3$ positive matrices, equality forces commutativity, while $n=4$ is the smallest dimension in which noncommuting equality examples exist; we further show this threshold is intrinsically about tensor-irreducibility: equality still forces commutativity in any total ambient dimension, provided the operators factor as pure Kronecker products $X=X_1\otimes\cdots\otimes X_k$ and $Y=Y_1\otimes\cdots\otimes Y_k$ with every factor $X_i,Y_i$ of dimension at most $3$. We derive contraction properties, sum inequalities, trace estimates for density operators, a Lipschitz-continuity property (including at singular limits), and limiting behaviour for Ces\`aro means and infinite products. As a worked example class, we specialize to density operators: we obtain an exact closed-form bijection between OMC and von Neumann entropy for qubits, prove that no such bijection exists for $d\ge3$, and establish explicit upper and lower envelopes for the entropy at fixed condition number, showing that the admissible entropy range $[S_{\min}(\kappa), S_{\max}(\kappa)]$ satisfies $\ln 2 \le S_{\max}(\kappa) < \ln 3$ for every $\kappa>1$. We also show that the qubit case recovers, as an operator lift, the classical equivalence between Michelson contrast and Jensen--Shannon divergence established by Bruni, Rossi, and Vitulano \cite{bruni2012}, and we give an explicit counterexample showing this equivalence provably fails to extend past $d=2$, as a direct consequence of the non-bijectivity above.
\end{abstract}

\noindent\textbf{2020 Mathematics Subject Classification.} Primary 47A10, 47B65, 47A30. Secondary 47B15, 53B20, 53C22, 81P45, 81Q10.

\noindent\textbf{Keywords.} Michelson contrast; logarithmic contrast; operator condition number; positive operators; Thompson metric; projective metric geometry; operator geometry; matrix inequalities; submultiplicativity; hyperbolic geometry; rapidity parametrisation; tensor products; von Neumann entropy.


\section{Introduction}

The classical Michelson contrast for two positive levels $0<m\le M$ is defined as
\[
\MC(M,m) = \frac{M-m}{M+m}.
\]
It satisfies the elementary identity
\[
\operatorname{arctanh}(\MC(M,m)) = \tfrac12\ln\frac{M}{m}.
\]
Thus the Michelson contrast may be viewed as a bounded parametrisation of the logarithmic ratio $M/m$. In logarithmic image processing, logarithmic representations of contrast have been developed; in particular, the logarithmic additive contrast (LAC) is closely related to the logarithmic ratio and the multiplicative structure of intensity levels. The works of Jourlin, Pinoli, and collaborators \cite{jourlin1988,jourlin1989} provide a classical foundation for this approach.

In a recent paper \cite{abed2024}, we introduced an operator analogue of the Michelson contrast for positive invertible operators $A$ on finite-dimensional spaces. The present work generalizes this construction to operators on arbitrary complex Hilbert spaces. For any invertible positive operator $A\in B(H)_{++}$ (where $B(H)$ denotes bounded operators on a Hilbert space $H$ of any dimension), define
\[
\Delta(A) = \frac{\|A\|-\|A^{-1}\|^{-1}}{\|A\|+\|A^{-1}\|^{-1}} = \frac{\kappa(A)-1}{\kappa(A)+1}, \qquad \kappa(A):=\|A\|\|A^{-1}\|.
\]
This operator Michelson contrast extends the classical formula by replacing the two scalar levels with the spectral supremum and infimum of the operator. It depends only on the condition number $\kappa(A)=\sup\sigma(A)/\inf\sigma(A)$, which is intrinsically dimension-independent.

The present paper investigates the logarithmic coordinate naturally associated with this operator Michelson contrast, in both finite and infinite dimensions:
\[
\QLC(A) := \operatorname{arctanh}(\Delta(A)) = \tfrac12\ln\kappa(A).
\]
We call this quantity the \emph{Operator Logarithmic Contrast} (OLC); the name is proposed in this work and is not an established term in the literature. It is the natural logarithmic coordinate for the operator Michelson contrast of \cite{abed2024}, and requires no physical interpretation to be well defined: it is a functional on invertible operators, studied here purely as such. Where the paper later specializes to density operators (Section \ref{sec:applications}), that specialization is presented as a worked example class -- a convenient, well-studied family of positive operators on which the general theory can be illustrated concretely -- rather than as the paper's central claim.

This logarithmic parametrisation reveals several new properties, which we develop first in complete generality for any Hilbert space, then specialize to finite dimensions where finer structure becomes visible. First, unlike the classical logarithmic ratio which is additive under multiplication of scalar ratios, the operator version satisfies only a subadditivity inequality, valid universally:
\[
\QLC(AB) \le \QLC(A) + \QLC(B),
\]
which follows from the submultiplicativity of the condition number and holds for all invertible operators in any Hilbert space. This inequality extends to limit-theoretic contexts: for convergent infinite products and Ces\`aro means of operator sequences, analogous inequalities hold by continuity and weak-operator topology arguments (Remark 4.2, Section \ref{sec:inf-dim}). Second, we prove a geometric identification: for positive invertible operators, $\QLC(A) = \delta_T([I],[A])$, where $\delta_T$ is the projective Thompson metric on the positive cone; this immediately yields a Lipschitz-continuity property for OLC that we make explicit below. Third, the equality case in the subadditivity inequality leads to an alignment condition for the extremal eigenspaces: $n=4$ is the smallest dimension in which noncommuting equality cases can occur, and we show this threshold is really a statement about tensor-irreducibility rather than raw ambient dimension. Applied to quantum information, these results give an exact closed-form relation between OMC and von Neumann entropy for qubits, and a proof, with explicit and now fully rigorous bounds, that no such relation exists in higher dimension.

The paper is organised as follows. Section 2 fixes notation and definitions, extending OLC to all invertible operators via the polar decomposition. Section 3 recalls the Thompson metric, proves the geometric identification theorem, and derives the continuity properties of OLC. Section 4 contains the linear-algebraic core: subadditivity, the equality case and dimensional threshold, tensor-product behaviour, power laws, contraction, sum inequalities, trace estimates, and double-sided products. Section 5 works out density operators as an example class. Section 6 discusses structural interpretations, open questions, and concludes.


\section{Definitions and Notation}

Let $H$ be a complex Hilbert space and $B(H)$ the algebra of all bounded operators. Let $B(H)_{++}$ and $B(H)_+$ denote the positive invertible and positive (including singular) operators, respectively.

For invertible $A\in B(H)$, the condition number is $\kappa(A):=\|A\|\|A^{-1}\|$. We recall the standard facts (see \cite{dixmier1969}):
\[
\kappa(A)\ge1,\qquad \kappa(AB)\le\kappa(A)\kappa(B),
\]
with equality in the first iff $A$ is a scalar multiple of a unitary operator. For $A\in B(H)_{++}$, $\kappa(A) = \sup\sigma(A)/\inf\sigma(A)$.

\subsection{The Operator Michelson Contrast}
In \cite{abed2024}, the operator Michelson contrast for positive operators was defined as
\[
\Delta(A) = \begin{cases} \dfrac{\|A\|-\|A^{-1}\|^{-1}}{\|A\|+\|A^{-1}\|^{-1}}, & A\in B(H)_{++},\\[2mm] 1, & A \text{ singular}. \end{cases}
\]
For invertible positive $A$, this equals $(\kappa(A)-1)/(\kappa(A)+1)$.

\subsection{Logarithmic Parametrisation}

\begin{definition}\label{def:qmc}
For any invertible $A\in B(H)$, define the \emph{Operator Michelson Contrast} (OMC) as the functional extension of $\Delta$:
\[
\QMC(A) := \frac{\kappa(A)-1}{\kappa(A)+1},
\]
and its logarithmic parametrisation, the \emph{Operator Logarithmic Contrast} (OLC), by
\[
\QLC(A) := \operatorname{arctanh}(\QMC(A)) = \tfrac12\ln\kappa(A).
\]
For singular nonzero operators, we may extend these quantities by the conventions $\QMC(A)=1$ and $\QLC(A)=+\infty$. The zero operator is excluded from this extended convention.
\end{definition}

For positive invertible $A$, this coincides with $\Delta(A)$ from \cite{abed2024}. By the geometric identification of Theorem \ref{thm:thompson} below, for positive invertible $A$, $\QLC(A)=\delta_T([I],[A])$.

\subsection{Extension to All Invertible Operators via Polar Decomposition}
Since $\kappa(A)=\kappa(|A|)$ for any invertible $A$ (polar decomposition $A=U|A|$), we have
\[
\QLC(A) = \tfrac12\ln\kappa(A) = \tfrac12\ln\kappa(|A|) = \QLC(|A|) = \delta_T([I],[|A|]).
\]
Thus the logarithmic contrast functional OLC on arbitrary invertible operators is obtained by pulling back the projective Thompson distance-to-identity functional through the absolute value map $A\mapsto|A|$. This is a natural extension of the functional, not of the metric itself, since the projective Thompson metric is originally defined only on the positive cone.

\begin{remark}\label{rem:polar}
This reduction is valid only for statements about $\kappa$, OMC, or OLC themselves; it does not transfer commutativity statements. If $X=\sigma_x$, $Y=\sigma_y$ are Pauli matrices, then $|X|=|Y|=I$ commute trivially, while $XY=i\sigma_z=-YX$ are maximally noncommuting. The dimensional-threshold result of Theorem \ref{thm:threshold} is intrinsically a statement about positive operators $X,Y$, using their own extremal eigenspaces $\Emax,\Emin$; it does not extend to general invertible operators via $|X|,|Y|$. See also Remark \ref{rem:entanglement}.
\end{remark}

Both quantities are invariant under positive scaling: for $\lambda>0$, $\QMC(\lambda A)=\QMC(A)$, $\QLC(\lambda A)=\QLC(A)$.


\section{Preliminaries: The Thompson Metric, Geometric Interpretation, and Continuity}

The Thompson metric on $B(H)_{++}$ (see \cite{thompson1963,corach1994}) is
\[
d_T(A,B) := \|\log(A^{-1/2}BA^{-1/2})\|.
\]
The projective version on rays $[A]=\{\lambda A:\lambda>0\}$ is
\[
\delta_T([A],[B]) := \inf_{\lambda>0} d_T(A,\lambda B).
\]
Further background on the Thompson metric and its relation to projective metrics can be found in \cite{lawson2001}.

\begin{theorem}\label{thm:thompson}
For any $A\in B(H)_{++}$,
\[
\delta_T([I],[A]) = \tfrac12\ln\kappa(A). \tag{1}
\]
\end{theorem}
\begin{proof}
Let $m=\inf\sigma(A)$, $M=\sup\sigma(A)$. Then $\kappa(A)=M/m$. For any $\lambda>0$,
\[
d_T(I,\lambda A) = \|\log(\lambda A)\| = \max\{|\log(\lambda m)|,|\log(\lambda M)|\}.
\]
Minimising over $\lambda>0$ gives $\lambda_*=(mM)^{-1/2}$ and the minimum equals $\tfrac12\ln(M/m)=\tfrac12\ln\kappa(A)$.
\end{proof}

Thus, for positive invertible operators, $\QLC(A)=\tfrac12\ln\kappa(A)$ is exactly the distance from the identity ray to the ray generated by $A$ in the projective Thompson metric.

\subsection{Continuity of OLC}

Theorem \ref{thm:thompson} has an immediate and useful consequence: since OLC \emph{is} a distance-to-a-fixed-point functional on a metric space, it inherits a Lipschitz bound for free from the triangle inequality. We record this, together with a companion statement showing that the singular convention $\QLC=+\infty$ of Definition \ref{def:qmc} is not an arbitrary choice but the actual limiting value of OLC as an invertible operator degenerates.

\begin{corollary}[1-Lipschitz continuity w.r.t.\ the Thompson metric]\label{cor:lipschitz}
For all $A,B\in B(H)_{++}$,
\[
|\QLC(A)-\QLC(B)| \le \delta_T([A],[B]).
\]
\end{corollary}
\begin{proof}
By Theorem \ref{thm:thompson}, $\QLC(A)=\delta_T([I],[A])$ and $\QLC(B)=\delta_T([I],[B])$. Since $\delta_T$ is a metric on rays, the triangle inequality gives $\delta_T([I],[A]) \le \delta_T([I],[B]) + \delta_T([B],[A])$ and symmetrically, so $|\delta_T([I],[A])-\delta_T([I],[B])| \le \delta_T([A],[B])$.
\end{proof}

In particular, on any subset of $B(H)_{++}$ on which $\delta_T$ is comparable to the operator norm (e.g.\ a norm-bounded set with condition numbers bounded away from the singular boundary), Corollary \ref{cor:lipschitz} yields ordinary norm-Lipschitz continuity of OLC.

\begin{proposition}[Continuity at singular limits]\label{prop:singular-cont}
Let $A_n\in B(H)$ be invertible with $A_n\to A$ in operator norm, where $A\neq0$ is singular (not necessarily positive). Then $\kappa(A_n)\to+\infty$, hence $\QLC(A_n)\to+\infty=\QLC(A)$ under the extended convention of Definition \ref{def:qmc}.
\end{proposition}
\begin{proof}
Norm convergence gives $\|A_n\|\to\|A\|$, and $\|A\|>0$ since $A\neq0$; in particular $\|A_n\|$ is bounded away from $0$ for large $n$. It remains to show $\|A_n^{-1}\|\to\infty$. Suppose not: then along a subsequence $\|A_n^{-1}\|\le M$ for some $M<\infty$. Since $A_n\to A$, for $n$ large enough $\|A_n-A\|<1/M$, and the standard Neumann-series perturbation bound for invertible operators shows that $A = A_n - (A_n-A) = A_n(I - A_n^{-1}(A_n-A))$ is invertible whenever $\|A_n^{-1}(A_n-A)\|\le\|A_n^{-1}\|\|A_n-A\|<1$, contradicting the singularity of $A$. Hence $\|A_n^{-1}\|\to\infty$ along every subsequence, so $\kappa(A_n)=\|A_n\|\|A_n^{-1}\|\to\infty$, and $\QLC(A_n)=\tfrac12\ln\kappa(A_n)\to+\infty$.
\end{proof}

Thus the extended convention $\QLC(A)=+\infty$ for nonzero singular $A$ is exactly the value forced by continuity from the invertible operators, not an independent stipulation.

\begin{remark}[An intermediate regime between finite dimensions and general $B(H)$]\label{rem:infinite-dim}
The dimensional-threshold results of Section \ref{sec:threshold} (Lemma \ref{lem:align} and Theorem \ref{thm:threshold}) are stated, and proved, only for finite-dimensional $X,Y\in M_n(\mathbb{C})_{++}$; the argument uses genuinely finite-dimensional facts (eigenspaces of finite multiplicity, orthogonal-complement invariance). A natural intermediate setting, between this finite-dimensional case and the fully general $B(H)_{++}$ for which only the metric identification of Theorem \ref{thm:thompson} is available, is that of positive invertible operators $X=I+K_X$, $Y=I+K_Y$ with $K_X,K_Y$ compact, or more generally operators whose spectrum accumulates only at points other than $\sup\sigma$ and $\inf\sigma$ (so that $\Emax$ and $\Emin$ are again finite-dimensional eigenspaces, as in the finite-dimensional proof). In that setting, the proof of Lemma \ref{lem:align} goes through essentially unchanged, since it only uses that vectors realising the operator norm exist and lie in a well-defined maximal/minimal eigenspace, which compactness of $K_X,K_Y$ guarantees at the extremal points of the spectrum. We do not carry out the corresponding case analysis of Theorem \ref{thm:threshold} in this setting, since the orthogonal-complement argument used there for $n=3$ requires some care when $\Emax,\Emin$ are infinite-dimensional or when the ``middle'' spectral subspace is infinite-dimensional; we record this as a natural direction for future work rather than a result of the present paper.
\end{remark}


\section{Core Linear-Algebraic Statements}

\subsection{Subadditivity and Composition Laws}

\begin{theorem}[Subadditivity in infinite and finite dimensions]\label{thm:subadd}
For any invertible $A,B\in B(H)$ (where $H$ is any complex Hilbert space, finite or infinite-dimensional),
\[
\QLC(AB) \le \QLC(A)+\QLC(B). \tag{2}
\]
Equivalently,
\[
\QMC(AB) \le \frac{\QMC(A)+\QMC(B)}{1+\QMC(A)\QMC(B)}. \tag{3}
\]
\end{theorem}
\begin{proof}
The proof is identical in finite and infinite dimensions: use the submultiplicativity $\kappa(AB)\le\kappa(A)\kappa(B)$, which holds for any bounded invertible operators. Taking logarithms gives (2); applying $\tanh$ yields (3).
\end{proof}

Thus OLC is subadditive rather than additive. Strict inequality can occur even for commuting operators if their extremal spectral directions are not aligned; see below.

\begin{remark}[Convergence and infinite products in limit-theoretic setting]
Theorem \ref{thm:subadd} immediately implies: if $A_1, A_2, A_3, \ldots\in B(H)_{++}$ is a sequence of invertible positive operators and the partial products $P_N:=A_1A_2\cdots A_N$ converge in operator norm to an invertible limit $P$, then
\[
\QLC(P) \le \liminf_{N\to\infty}\sum_{k=1}^N\QLC(A_k).
\]
If the series $\sum_{k=1}^\infty\QLC(A_k)$ converges, then $\QLC(P)\le\sum_{k=1}^\infty\QLC(A_k)$. This applies equally in infinite-dimensional Hilbert spaces: the operator-norm convergence and the continuity of OLC (Corollary \ref{cor:lipschitz}, Proposition \ref{prop:singular-cont}) require no assumption on the dimension of $H$. 

Similarly, for Ces\`aro means $S_N=\tfrac1N\sum_{k=1}^N A_k$ of a sequence of positive invertible operators, the sum inequality (9) gives $\QLC(S_N)\le\max_k\QLC(A_k)$, and if $S_N$ converges to $A$ invertible in operator norm, then $\QLC(A)\le\sup_k\QLC(A_k)$ by continuity.
\end{remark}

\subsection{Equality Case and Dimensional Threshold}\label{sec:threshold}

For finite-dimensional positive matrices $X,Y\in M_n(\mathbb{C})_{++}$, equality in (2) is equivalent to
\[
\|XY\| = \|X\|\|Y\|, \qquad \|(XY)^{-1}\| = \|X^{-1}\|\|Y^{-1}\|. \tag{4}
\]

\begin{lemma}[Extremal subspace alignment via spectral projections]\label{lem:align}
Let $X,Y\in B(H)_{++}$ (with $H$ any Hilbert space) satisfy (4). Let $E_{\max}(X)$ denote the spectral subspace of $X$ corresponding to its largest eigenvalue $\lambda_{\max}(X):=\|X\|$ (defined as the range of the spectral projection $\mathbf{1}_{[\lambda_{\max}(X),\infty)}(X)$), and $E_{\min}(X)$ the spectral subspace at its smallest eigenvalue $\lambda_{\min}(X)$ (the range of $\mathbf{1}_{\{\lambda_{\min}(X)\}}(X)$). Then
\[
E_{\max}(X)\cap E_{\max}(Y)\neq\{0\}, \qquad E_{\min}(X)\cap E_{\min}(Y)\neq\{0\}.
\]
The same statement holds if $E_{\max}, E_{\min}$ denote the extremal generalised eigenspaces (i.e., kernels of $(X-\lambda_{\max}I)^k$ and $(X-\lambda_{\min}I)^k$ for sufficiently large $k$), relevant when multiplicities are greater than one.
\end{lemma}
\begin{proof}
From $\|XY\|=\|X\|\|Y\|$, the extremal norm conditions (Fact in proof of Theorem \ref{thm:subadd}) applied to the product $XY$ yield: there exists a unit vector $w$ (or more generally, a sequence of unit vectors for infinite rank) with $\|XYw\|=\|XY\|\to\|X\|\|Y\|$ in the limit. The chain of norm inequalities $\|XYw\|\le\|X\|\|Yw\|\le\|X\|\|Y\|$ forces both to be equalities, hence $\|Yw\|=\|Y\|$ and $\|X(Yw)\|=\|X\|\|Yw\|$. This means $w$ is a maximal-norm vector for $Y$ and $Yw$ is a maximal-norm vector for $X$. By the definition of spectral subspaces (equivalently, by the compactness of $E_{\max}(Y)$ when $E_{\max}(Y)$ is finite-dimensional), $w\in E_{\max}(Y)$, and by invertibility of $Y$, $Yw\in E_{\max}(Y)$ as well. Moreover, $Yw\in E_{\max}(X)$. Thus $Yw\in E_{\max}(X)\cap E_{\max}(Y)$, which is nonzero. The minimal part follows by applying the argument to $X^{-1},Y^{-1}$.
\end{proof}

\begin{theorem}[Dimensional threshold for equality]\label{thm:threshold}
Let $X,Y\in M_n(\mathbb{C})_{++}$. If $\QLC(XY)=\QLC(X)+\QLC(Y)$, then $XY=YX$ for $n\le3$. Moreover, $n=4$ is the smallest dimension in which noncommuting equality cases can occur.
\end{theorem}
\begin{proof}
For $n=1$ trivial.

\emph{Case $n=2$.} Let $X=\operatorname{diag}(\lambda_1,\lambda_2)$, $\lambda_1\ge\lambda_2>0$, and $Y=\begin{pmatrix}a&b\\\bar b&c\end{pmatrix}>0$. If $\lambda_1=\lambda_2$, $X$ is scalar and commutes. Assume $\lambda_1>\lambda_2$. By Lemma \ref{lem:align}, there exists $u\in\Emax(X)\cap\Emax(Y)$. Since $\Emax(X)=\operatorname{span}\{e_1\}$, $u$ is proportional to $e_1$, so $Ye_1=\mu e_1$, giving $b=0$. Similarly, the minimal alignment gives $\Emin(X)\cap\Emin(Y)\neq\{0\}$, so $e_2$ is an eigenvector of $Y$. Hence $Y$ is diagonal and commutes with $X$.

\emph{Case $n=3$.} Let $X=\operatorname{diag}(\lambda_1,\lambda_2,\lambda_3)$, $\lambda_1\ge\lambda_2\ge\lambda_3>0$, and $Y>0$.

\emph{Subcase $\lambda_1>\lambda_2>\lambda_3$.} $\Emax(X)=\operatorname{span}\{e_1\}$ is one-dimensional; alignment gives $Ye_1=\mu_1e_1$. Similarly $\Emin(X)=\operatorname{span}\{e_3\}$ gives $Ye_3=\mu_3e_3$. Since $Y$ is self-adjoint, $\operatorname{span}\{e_1\}$ and $\operatorname{span}\{e_3\}$ are $Y$-invariant, hence so is their orthogonal complement $\operatorname{span}\{e_2\}$; thus $Y$ is diagonal in the basis $\{e_1,e_2,e_3\}$.

\emph{Subcase $\lambda_1=\lambda_2>\lambda_3$.} Let $E=\operatorname{span}\{e_1,e_2\}$, $F=\operatorname{span}\{e_3\}$. Minimal alignment gives $e_3\in\Emin(Y)$ (since $\Emin(X)=F$), so $Ye_3=\mu_3e_3$; thus $F$ is $Y$-invariant, hence so is $E=F^\perp$. On $E$, $X$ is scalar ($\lambda_1I_E$), so $Y$ commutes with $X$ on $E$; on $F$, $X$ is scalar ($\lambda_3I_F$), so $Y$ commutes there too. Therefore $XY=YX$.

\emph{Subcase $\lambda_1>\lambda_2=\lambda_3$.} Symmetric argument using maximal alignment.

Thus for $n\le3$, equality forces commutativity.

\emph{Noncommuting example for $n=4$.} Let
\[
X=\begin{pmatrix}4&0&0&0\\0&3&0&0\\0&0&1&0\\0&0&0&1\end{pmatrix},\qquad
Y=\begin{pmatrix}4&0&0&0\\0&2&1&0\\0&1&2&0\\0&0&0&1\end{pmatrix}.
\]
Then $\kappa(X)=4$, and $\sigma(Y)=\{4,3,1,1\}$, so $\kappa(Y)=4$. We have
\[
XY=\begin{pmatrix}16&0&0&0\\0&6&3&0\\0&1&2&0\\0&0&0&1\end{pmatrix}.
\]
Let $B=\begin{pmatrix}6&3\\1&2\end{pmatrix}$. Then $\|B\|\le\|B\|_F=\sqrt{50}<16$, and $B^{-1}=\tfrac19\begin{pmatrix}2&-3\\-1&6\end{pmatrix}$, $\|B^{-1}\|\le\|B^{-1}\|_F=\sqrt{50}/9<1$. Since $XY$ has the scalar blocks $16$ and $1$, we get $\|XY\|=16$ and $\|(XY)^{-1}\|=1$. Hence $\kappa(XY)=16=4\cdot4$, so equality in (2) holds. However, the blocks on $e_2,e_3$ satisfy
\[
\operatorname{diag}(3,1)\begin{pmatrix}2&1\\1&2\end{pmatrix} \neq \begin{pmatrix}2&1\\1&2\end{pmatrix}\operatorname{diag}(3,1),
\]
so $X$ and $Y$ do not commute. This establishes that noncommuting equality cases exist for $n=4$, and hence for all $n\ge4$ by taking direct sums with identity blocks.
\end{proof}

\begin{remark}[The intermediate infinite-dimensional regime]
Theorem \ref{thm:threshold} is stated and proved for finite-dimensional positive matrices $X,Y\in M_n(\mathbb{C})_{++}$. The argument uses the finiteness of eigenspace dimensions and the orthogonal-complement structure of $\mathbb{C}^n$. 

However, for certain infinite-dimensional operators, the result extends via a spectral-measure perspective. Suppose $X,Y\in B(H)_{++}$ satisfy the equality condition (4), and suppose further that:
\begin{enumerate}
\item The spectrum of $X$ consists of isolated points (i.e., point spectrum) with finite multiplicity at each point.
\item The extremal eigenvalues $\lambda_{\max}(X),\lambda_{\min}(X)$ are isolated (not accumulation points of other spectrum), with finite-dimensional eigenspaces $E_{\max}(X),E_{\min}(X)$.
\item The ``middle'' part of the spectrum (eigenvalues strictly between the extrema) either is empty, or is itself finite-dimensional.
\end{enumerate}

Under these conditions, the restriction of $X,Y$ to the finite-dimensional invariant subspace spanned by the extremal and intermediate eigenspaces of $X$ (or $Y$) reduces to the finite-dimensional case, and the commutativity conclusion of Theorem \ref{thm:threshold} follows. Examples include: 
- Compact perturbations $X=I+K$ with $K$ compact of finite rank.
- Countable multiplicity with finite total dimension ($\dim E_{\max}\oplus E_{\text{middle}}\oplus E_{\min}<\infty$).
- Operators defined by multiplication on $L^2(\mathbb{R})$ with support on a finite interval, whose spectrum is discrete.

For fully general $B(H)_{++}$, the threshold result and commutativity conclusions do not apply without additional hypotheses, since the extremal subspaces may be infinite-dimensional and the case analysis of Theorem \ref{thm:threshold} for intermediate multiplicity becomes subtle. We leave a complete characterization of the infinite-dimensional threshold structure as a natural direction for future work (see also Open Questions).
\end{remark}

\begin{remark}
The alignment phenomenon appearing in the equality case is reminiscent of the algebraic structures associated with K-loops and gyrovector spaces on the positive cone (see \cite{beneduci2014,abe2015,park2015,kim2016}), but we do not claim that our equality condition is equivalent to vanishing gyration in general; this remains an open problem. Our matrix proof for $n\le3$ gives a self-contained verification in low dimensions, and the alignment principle (Lemma \ref{lem:align}) extends to the intermediate infinite-dimensional regime.
\end{remark}

\subsection{Tensor Products and Composite Systems}

\begin{proposition}[Tensor product additivity]\label{prop:tensor-add}
For invertible $A,B$, the singular values of $A\otimes B$ are products of the singular values of $A$ and of $B$, hence
\[
\kappa(A\otimes B)=\kappa(A)\kappa(B), \qquad \QLC(A\otimes B)=\QLC(A)+\QLC(B).
\]
\end{proposition}
\begin{proof}
Immediate from the singular value decomposition of $A\otimes B$.
\end{proof}

\begin{proposition}[Local vs. global contrast]\label{prop:local-global}
Let $Z\in B(H_1\otimes H_2)$ be invertible, $\dim H_2=d_2<\infty$. Then
\[
\QLC(Z) \ge \max\{\QLC(\operatorname{Tr}_{H_2}|Z|), \QLC(\operatorname{Tr}_{H_1}|Z|)\}.
\]
\end{proposition}
\begin{proof}
Write $Z>0$ (the general case follows by applying the argument to $|Z|$; see Remark \ref{rem:polar} for why this reduction is valid here but not for commutativity statements). The operator inequalities $\lambda_{\min}(Z)I\le Z\le\lambda_{\max}(Z)I$ are preserved under the positive partial trace. Since $\operatorname{Tr}_{H_2}I=d_2I_1$,
\[
d_2\lambda_{\min}(Z)I_1 \le \operatorname{Tr}_{H_2}Z \le d_2\lambda_{\max}(Z)I_1,
\]
so $\lambda_{\min}(\operatorname{Tr}_{H_2}Z)\ge d_2\lambda_{\min}(Z)$ and $\lambda_{\max}(\operatorname{Tr}_{H_2}Z)\le d_2\lambda_{\max}(Z)$. Hence
\[
\kappa(\operatorname{Tr}_{H_2}Z) \le \frac{d_2\lambda_{\max}(Z)}{d_2\lambda_{\min}(Z)} = \kappa(Z),
\]
and taking logarithms gives $\QLC(\operatorname{Tr}_{H_2}Z)\le\QLC(Z)$. The same argument applies to the other partial trace.
\end{proof}

\begin{remark}
No general upper bound exists: the counterexample $Z_\varepsilon = I\otimes I + \varepsilon(X\otimes X+Y\otimes Y)$ on $\mathbb{C}^2\otimes\mathbb{C}^2$ gives positive $\QLC(Z_\varepsilon)$ while both partial traces are proportional to the identity ($\QLC=0$).
\end{remark}

\begin{proposition}[Chain products]
For a chain $H=\bigotimes_{i=1}^N H_i$ and nearest-neighbour invertible operators $A_{i,i+1}$,
\[
\QLC(A_{12}A_{23}\cdots A_{N-1,N}) \le \sum_{i=1}^{N-1}\QLC(A_{i,i+1}).
\]
Equality requires extremal subspace alignment at each step.
\end{proposition}
\begin{proof}
Repeated application of Theorem \ref{thm:subadd}; equality conditions follow from Lemma \ref{lem:align}.
\end{proof}

\begin{proposition}[Tensor powers]
$\QLC(A^{\otimes N})=N\,\QLC(A)$.
\end{proposition}
\begin{proof}
Immediate from Proposition \ref{prop:tensor-add}.
\end{proof}

We now show that the dimensional threshold of Theorem \ref{thm:threshold} interacts sharply with tensor-product additivity, giving a criterion independent of total ambient dimension. The key technical step is the following elementary but easily mis-stated fact about products of bounded positive reals, which we now prove carefully rather than by assertion.

\begin{lemma}[Tightness of a bounded product]\label{lem:tight-product}
Let $a_1,\dots,a_k,b_1,\dots,b_k$ be positive real numbers with $a_i\le b_i$ for every $i=1,\dots,k$, and suppose
\[
\prod_{i=1}^k a_i = \prod_{i=1}^k b_i.
\]
Then $a_i=b_i$ for every $i=1,\dots,k$.
\end{lemma}
\begin{proof}
We induct on $k$. The case $k=1$ is immediate. For the base case $k=2$: since $a_1\le b_1$, $a_2\le b_2$, and all four numbers are positive,
\[
a_1a_2 \le a_1b_2 \le b_1b_2,
\]
and by hypothesis $a_1a_2=b_1b_2$, so both inequalities are equalities. From $a_1b_2=b_1b_2$ and $b_2>0$ we get $a_1=b_1$. From $a_1a_2=a_1b_2$ and $a_1=b_1>0$ we get $a_2=b_2$.

Now assume the lemma holds for $k-1$ factors, and let $a_1,\dots,a_k,b_1,\dots,b_k>0$ satisfy $a_i\le b_i$ for all $i$ and $\prod_{i=1}^k a_i=\prod_{i=1}^k b_i$. Set $A:=\prod_{i=1}^{k-1}a_i$ and $B:=\prod_{i=1}^{k-1}b_i$. Since each $a_i\le b_i$ and all terms are positive, $A\le B$ (a product of positive numbers is monotone in each factor). The hypothesis reads $Aa_k=Bb_k$, with $A\le B$ and $a_k\le b_k$, both pairs positive; applying the $k=2$ base case to the pair $(A,a_k)$ and $(B,b_k)$ gives $A=B$ and $a_k=b_k$. Having established $A=B$, i.e.\ $\prod_{i=1}^{k-1}a_i=\prod_{i=1}^{k-1}b_i$ with $a_i\le b_i$ for $i=1,\dots,k-1$, the induction hypothesis gives $a_i=b_i$ for $i=1,\dots,k-1$. Together with $a_k=b_k$, this proves the claim for $k$ factors.
\end{proof}

\begin{proposition}[Equality transfer under pure tensor products]\label{prop:tensor-transfer}
Let $X=X_1\otimes\cdots\otimes X_k$ and $Y=Y_1\otimes\cdots\otimes Y_k$, with $X_i,Y_i\in M_{n_i}(\mathbb{C})_{++}$. If equality holds in the subadditivity law,
\[
\QLC(XY) = \QLC(X)+\QLC(Y),
\]
then equality holds factorwise: $\kappa(X_iY_i)=\kappa(X_i)\kappa(Y_i)$ for every $i=1,\dots,k$. In particular, if $n_i\le3$ for every $i$, then $X_iY_i=Y_iX_i$ for every $i$, and consequently $XY=YX$.
\end{proposition}
\begin{proof}
By Proposition \ref{prop:tensor-add}, $\kappa(X)=\prod_i\kappa(X_i)$, $\kappa(Y)=\prod_i\kappa(Y_i)$, and since $XY=\bigotimes_i X_iY_i$, also $\kappa(XY)=\prod_i\kappa(X_iY_i)$. The hypothesis $\kappa(XY)=\kappa(X)\kappa(Y)$ thus reads
\[
\prod_{i=1}^k \kappa(X_iY_i) = \prod_{i=1}^k \big[\kappa(X_i)\kappa(Y_i)\big],
\]
where each factor on the left is bounded above by the corresponding factor on the right, $\kappa(X_iY_i)\le\kappa(X_i)\kappa(Y_i)$, by ordinary submultiplicativity of $\kappa$. Applying Lemma \ref{lem:tight-product} with $a_i=\kappa(X_iY_i)$ and $b_i=\kappa(X_i)\kappa(Y_i)$ gives $\kappa(X_iY_i)=\kappa(X_i)\kappa(Y_i)$ for each $i$. When $n_i\le3$, Theorem \ref{thm:threshold} converts this into $X_iY_i=Y_iX_i$; since factors of a Kronecker product commute exactly when each corresponding pair of tensor legs commutes, $XY=YX$ follows.
\end{proof}

\begin{remark}
Proposition \ref{prop:tensor-transfer} shows that the threshold $n\le3$ in Theorem \ref{thm:threshold} is a statement about irreducibility rather than raw ambient dimension: equality still forces commutativity in arbitrarily large $n$, provided $X,Y$ factor as pure tensor products of blocks of size at most $3$. The noncommuting equality-saturating example of Theorem \ref{thm:threshold} therefore cannot be reproduced, at any total dimension, by padding with further tensor factors of size $\le3$; its irreducible $4\times4$ block is essential.
\end{remark}

\subsection{Power Laws and Interpolation}
For $A\in B(H)_{++}$ and $\alpha\in\mathbb{R}$,
\[
\QLC(A^\alpha) = |\alpha|\,\QLC(A). \tag{5}
\]
Indeed, $\kappa(A^\alpha)=\kappa(A)^{|\alpha|}$.

For $X,Y\in B(H)_{++}$ and $\alpha\in[0,1]$, define $M_\alpha=X^\alpha Y^{1-\alpha}$. Although $M_\alpha$ need not be positive when $X,Y$ do not commute, it is invertible, and $\QLC(M_\alpha)=\QLC(|M_\alpha|)$ by the polar-decomposition extension. Then from subadditivity,
\[
\QLC(|M_\alpha|) \le \alpha\,\QLC(X) + (1-\alpha)\,\QLC(Y). \tag{6}
\]

\subsection{Contraction and Spectral Range Compression}

The logarithmic contrast admits two distinct types of contraction principles.

\paragraph{Contraction from projective Thompson non-expansiveness.} Let $F:B(H)_{++}\to B(H)_{++}$ be non-expansive w.r.t.\ the projective Thompson metric, $\delta_T([F(A)],[F(B)])\le\delta_T([A],[B])$, and suppose $[F(I)]=[I]$. Then, using (1),
\[
\QLC(F(A)) = \delta_T([I],[F(A)]) = \delta_T([F(I)],[F(A)]) \le \delta_T([I],[A]) = \QLC(A).
\]
Thus OLC is contracted by such maps. Since $\QMC=\tanh(\QLC)$, the same holds for OMC.

\paragraph{Spectral range compression.} Suppose $F$ satisfies $M_{F(A)}\le M_A$, $m_{F(A)}\ge m_A$, where $M_A=\|A\|$, $m_A=\|A^{-1}\|^{-1}$. Then $M_{F(A)}/m_{F(A)}\le M_A/m_A$. Since $\QMC(A)=1-2/(M_A/m_A+1)$, $\QLC(A)=\tfrac12\ln(M_A/m_A)$, and both functions are increasing on $[1,\infty)$, we obtain $\QMC(F(A))\le\QMC(A)$, $\QLC(F(A))\le\QLC(A)$.

Unlike the preceding argument, this conclusion does not use the Thompson metric or any non-expansiveness property of $F$; it follows solely from the compression of the spectral range.

\paragraph{Contrast expansion.} The same spectral reasoning gives the converse criterion: if $M_{G(A)}\ge M_A$, $m_{G(A)}\le m_A$, then $\kappa(G(A))\ge\kappa(A)$, and consequently $\QLC(G(A))\ge\QLC(A)$, $\QMC(G(A))\ge\QMC(A)$.

\subsection{Lower Bound via Deviation from Scalar Operators}
For any $A>0$,
\[
\QLC(A) \ge \tfrac12\ln\Big(1+\Big\|I-\frac{A}{\|A\|}\Big\|\Big). \tag{7}
\]
The proof is elementary using $m=\lambda_{\min}(A)$, $M=\lambda_{\max}(A)$; equality holds iff $A$ is scalar.

\subsection{Boundary Behaviour}
If $A_n>0$ satisfies $\sup\sigma(A_n)/\inf\sigma(A_n)\to\infty$, then $\QLC(A_n)\to\infty$. In particular, if $\inf\sigma(A_n)\to0$ while $\sup\sigma(A_n)$ remains bounded away from zero, $\QLC(A_n)\to\infty$ (consistent with Proposition \ref{prop:singular-cont}). If both extremal values tend to zero at comparable rates, the contrast may remain finite.

\subsection{Sum Inequality and Limiting Behaviour}
For positive operators (see \cite{abed2024}),
\[
\QLC(A+B) \le \max\{\QLC(A),\QLC(B)\}. \tag{8}
\]
By iteration, for $A_1,\dots,A_N>0$,
\[
\QLC\Big(\sum_{k=1}^N A_k\Big) \le \max_{1\le k\le N}\QLC(A_k). \tag{9}
\]
For Ces\`aro means $S_N=\tfrac1N\sum_{k=1}^N A_k$, $\QLC(S_N)\le\max_{1\le k\le N}\QLC(A_k)$, and if $S_N\to A$ invertible, then
\[
\QLC(A) \le \limsup_{N\to\infty}\QLC(S_N) \le \sup_{k\ge1}\QLC(A_k). \tag{10}
\]
For infinite products $P_N=A_1\cdots A_N$,
\[
\QLC(P_N) \le \sum_{k=1}^N \QLC(A_k), \tag{11}
\]
and if the series converges and $P_N\to P$ invertible, then
\[
\QLC(P) \le \sum_{k=1}^\infty \QLC(A_k). \tag{12}
\]
Even in the commutative case, equality requires alignment of extremal spectral directions.

\subsection{Trace Inequality for Density Operators (Finite Dimensional)}
Let $\dim H=n<\infty$ and $A>0$ with $\operatorname{Tr}(A)=1$. Let $\lambda_{\min}(A)$ be its smallest eigenvalue. Then
\[
\QLC(A) \le -\tfrac12\ln\lambda_{\min}(A). \tag{13}
\]
\begin{proof}
Let $\lambda_{\max}(A)=\|A\|$. Since $\operatorname{Tr}(A)=1$, $\lambda_{\max}(A)\le1$. Hence $\kappa(A)=\lambda_{\max}(A)/\lambda_{\min}(A)\le1/\lambda_{\min}(A)$. Taking logarithms and dividing by 2 yields (13).
\end{proof}
This bound is asymptotically sharp as $\lambda_{\min}\to0$; equality cannot be attained for strictly positive density matrices.

\subsection{Double-Sided Products}
A natural class of operators is given by double-sided products $Z=X^{1/2}YX^{1/2}$, where $X>0$ and $Y$ is invertible. $Z$ is invertible; if $Y=Y^*$, $Z$ is self-adjoint; if $Y>0$, $Z>0$.

The general submultiplicativity of the condition number gives
\[
\kappa(X^{1/2}YX^{1/2}) \le \kappa(X^{1/2})\kappa(Y)\kappa(X^{1/2}) = \kappa(X)\kappa(Y).
\]
Consequently,
\[
\QLC(X^{1/2}YX^{1/2}) \le \QLC(X)+\QLC(Y). \tag{14}
\]
When $Y>0$, setting $T=X^{1/2}Y^{1/2}$ gives $TT^*=X^{1/2}YX^{1/2}=Z$, and since the eigenvalues of $TT^*$ are the squares of the singular values of $T$, $\kappa(Z)=\kappa(T)^2$. Hence $\kappa(Z)=\kappa(X^{1/2}Y^{1/2})^2\le\kappa(X)\kappa(Y)$, recovering (14). Thus positivity of $Y$ is not required for the inequality itself; it only provides the additional factorisation through $Y^{1/2}$.


\section{A Worked Example Class: Density Operators}\label{sec:applications}

Density operators are a natural, well-studied family of positive operators, and they let us make the preceding general theory concrete: OLC coincides with a familiar quantity (a spectral gap statistic of $\rho$), von Neumann entropy gives a second, independently well-understood functional to compare it against, and the qubit/qutrit cases are small enough to compute exactly. None of what follows depends on quantum-mechanical structure specific to density operators (tensor products with a physical subsystem interpretation, measurement, dynamics); the results hold for any positive invertible operator of trace $1$, and we use standard quantum-information terminology (state, entropy, divergence) only because it is the terminology this example family is normally discussed in.

\subsection{Density Matrices}

\begin{proposition}
For a density matrix $\rho$ (trace 1, positive), if $\rho$ is invertible, then
\[
\QLC(\rho) = \tfrac12\ln\frac{\lambda_{\max}(\rho)}{\lambda_{\min}(\rho)}.
\]
It vanishes for the maximally mixed state $\rho=I/d$ and tends to $+\infty$ as $\rho$ approaches a pure state. Pure states have $\QLC=+\infty$ by our singular convention.
\end{proposition}
\begin{proof}
Immediate from Definition \ref{def:qmc}.
\end{proof}

\begin{remark}[Color images and the $d=3$ appearance]
In classical logarithmic image processing (LIP) and the closely related Color LIP, a color image $\mathbf{I}=(R,G,B)$ (red, green, blue intensities) is represented in log-space as $\mathbf{H}=(\log R,\log G,\log B)$. The ``intensity matrix'' or ``channel dominance'' matrix $M_\mathbf{I}:=\operatorname{diag}(R,G,B)$ is a $3\times3$ positive diagonal operator with eigenvalues $(R,G,B)$, and its condition number $\kappa(M_\mathbf{I})=\max(R,G,B)/\min(R,G,B)$ measures the spread of color intensities. The OLC of this matrix is $\tfrac12\ln\kappa(M_\mathbf{I})=\QLC(M_\mathbf{I})$, which is exactly the log-contrast of the color image in the sense of Jourlin--Pinoli's original contrast theory. The fact that our paper independently identifies $d=3$ (dimension 3) as the first dimension in which the equality case of the subadditivity law admits noncommuting examples (Theorem \ref{thm:threshold}) is thus not coincidental: the dimensional threshold is a deep structure of the positive cone for matrices of all dimensions, and the appearance of $d=3$ in both color imaging (three channels) and in our abstract theorem reflects a real underlying geometry. We do not develop this connection further here, but it suggests that the projective-Thompson-metric viewpoint and the spectral subadditivity analysis may have natural applications in multi-channel image processing beyond classical Michelson contrast.
\end{remark}

\subsection{Relation to von Neumann Entropy}

\begin{proposition}\label{prop:entropy-bound}
Let $S(\rho)=-\sum_i\lambda_i\ln\lambda_i$ be the von Neumann entropy. Then for a density matrix $\rho$ on a $d$-dimensional space,
\[
\QLC(\rho) \ge \tfrac12(S(\rho)-\ln d). \tag{15}
\]
\end{proposition}
\begin{proof}
Since $\lambda_{\max}\ge1/d$ and $S\le-\ln\lambda_{\min}$ (as $-\ln\lambda_{\min}$ is the maximum of $-\ln\lambda_i$), we have $\lambda_{\min}\le e^{-S}$. Thus $\lambda_{\max}/\lambda_{\min}\ge e^S/d$, yielding the inequality.
\end{proof}

\begin{proposition}[Exact qubit relation]\label{prop:qubit}
Let $d=2$ and write $x=\QMC(\rho)\in[0,1)$. Then
\[
S(\rho) = \ln2 - \tfrac12\big[(1+x)\ln(1+x)+(1-x)\ln(1-x)\big]. \tag{16}
\]
Equivalently, in terms of $\eta=\QLC(\rho)=\operatorname{arctanh}(x)$,
\[
S(\rho) = \ln2 - \tfrac12\big[(1+\tanh\eta)\ln(1+\tanh\eta)+(1-\tanh\eta)\ln(1-\tanh\eta)\big].
\]
The map $x\mapsto S$ is a strictly decreasing bijection $[0,1)\to(0,\ln2]$, with $S=\ln2$ at $x=0$ (maximally mixed) and $S\to0$ as $x\to1$ (pure states).
\end{proposition}
\begin{proof}
Write the eigenvalues of $\rho$ as $p\ge1-p$, so $\kappa(\rho)=p/(1-p)$ and
\[
x = \QMC(\rho) = \frac{\kappa-1}{\kappa+1} = \frac{p-(1-p)}{p+(1-p)} = 2p-1,
\]
hence $p=\tfrac{1+x}2$, $1-p=\tfrac{1-x}2$. Substituting into $S(\rho)=-p\ln p-(1-p)\ln(1-p)$ and collecting the $\ln2$ terms yields (16). Monotonicity follows since $dS/dx=\tfrac12\ln\tfrac{1-x}{1+x}<0$ for $x\in(0,1)$, and the endpoint values are immediate.
\end{proof}

\begin{figure}[htbp]
\centering
\includegraphics[width=0.62\linewidth]{fig1_qmc_entropy.ps}
\caption{$S(\rho)$ as a function of $\QMC(\rho)$ for a qubit density matrix (Proposition \ref{prop:qubit}), linear scale.}
\end{figure}

\begin{figure}[htbp]
\centering
\includegraphics[width=0.62\linewidth]{fig2_qlc_entropy.ps}
\caption{$S(\rho)$ as a function of $\QLC(\rho)$ for a qubit density matrix, symmetric log-scale window $\QLC\in[10^{-2},10^2]$ centred on $\QLC=1$.}
\end{figure}

\begin{remark}\label{rem:qubit-remark}
Proposition \ref{prop:qubit} shows that for $d=2$, OMC (equivalently OLC) and $S$ carry exactly the same information, related by the strictly decreasing, strictly concave map (16) -- the binary entropy of the Bloch-sphere radial coordinate $p=(1+x)/2$. Note also that $dS/dx\to0$ as $x\to0$: $S$ is stationary to first order at the maximally mixed state, so OMC (and OLC) are more sensitive local coordinates than $S$ near maximal mixedness. For $d\ge3$ no bijection can hold, since $S$ depends on the full spectrum while OMC depends only on $\lambda_{\max},\lambda_{\min}$; only the one-sided bound (15) survives, as made precise in Section \ref{sec:nonbij}.
\end{remark}

\begin{remark}[No exact symmetry about $\eta=1$]\label{rem:no-symmetry}
The graphs of $S$ against $\QMC(\rho)$ and against $\eta:=\QLC(\rho)$ visually suggest a symmetry about $\eta=1$, but no such symmetry holds. If $S(\eta)+S(1/\eta)$ were constant, this would force $\eta\mapsto1/\eta$ to be an exact involution of the entropy graph; numerically, $S(e)+S(1/e)\approx0.02792+0.62980=0.65771$, which equals neither $\ln2\approx0.69315$ nor $2S(1)\approx0.73067$. The visual impression of near-symmetry is an artifact of $\tanh$'s transition from its linear regime to saturation occurring near $\eta\sim1$, not an exact property of $S$.
\end{remark}

\begin{corollary}[Asymptotics of the qubit entropy near maximally mixed and near pure states]\label{cor:entropy-asymp}
Let $\eta=\QLC(\rho)$ for a qubit density matrix $\rho$. Then:
\begin{enumerate}[label=(\alph*)]
\item As $\eta\to0$ (near the maximally mixed state),
\[
S(\eta) = \ln2 - \tfrac12\eta^2 + O(\eta^4);
\]
in particular $S$ is stationary to first order at $\eta=0$ (consistent with Remark \ref{rem:qubit-remark}), and there is no linear dependence of $S$ on $\eta$ near the maximally mixed state.
\item As $\eta\to\infty$ (near a pure state),
\[
S(\eta) = (1+2\eta)e^{-2\eta}\big(1+O(e^{-2\eta})\big),
\]
equivalently
\[
\ln S(\eta) = -2\eta + \ln(1+2\eta) + O(e^{-2\eta}).
\]
Thus $\ln S(\eta)-\ln(1+2\eta)$ is asymptotically linear in $\eta$ with slope $-2$; this linearisation already agrees with $\ln S(\eta)$ to within $5\times10^{-4}$ by $\eta\approx4$.
\end{enumerate}
Neither of these two asymptotic regimes is linear in $\eta$ globally, and by Remark \ref{rem:no-symmetry} they are not related to one another by any exact symmetry $\eta\leftrightarrow1/\eta$: (a) governs the small-$\eta$ plateau, (b) governs the large-$\eta$ tail, and the transition between the two occurs precisely in the intermediate region $\eta\sim1$ where $\tanh$ leaves its linear regime.
\end{corollary}
\begin{proof}
(a) Write $x=\tanh\eta$ and $f(x):=(1+x)\ln(1+x)+(1-x)\ln(1-x)$, so $S=\ln2-\tfrac12 f(x)$ by (16). Then $f(0)=0$, $f'(x)=\ln\frac{1+x}{1-x}$ so $f'(0)=0$, and $f''(x)=\frac2{1-x^2}$ so $f''(0)=2$; hence $f(x)=x^2+O(x^4)$ and $S=\ln2-\tfrac12x^2+O(x^4)$. Since $x=\tanh\eta=\eta+O(\eta^3)$, $x^2=\eta^2+O(\eta^4)$, giving (a).

(b) Let $\varepsilon:=1-x=1-\tanh\eta$, so $\varepsilon=\dfrac{2e^{-2\eta}}{1+e^{-2\eta}}=2e^{-2\eta}\big(1+O(e^{-2\eta})\big)$. Writing $p=(1+x)/2=1-\varepsilon/2$ and $q=1-p=\varepsilon/2$, the standard small-$q$ expansion of the binary entropy $S=-p\ln p-q\ln q$ gives $S=q(1-\ln q)+O(q^2)$, using $-(1-q)\ln(1-q)=q-\tfrac12q^2+O(q^3)$ together with the exact term $-q\ln q$. Substituting $q=\varepsilon/2$,
\[
S = \frac\varepsilon2\Big(1-\ln\frac\varepsilon2\Big)+O(\varepsilon^2) = \frac\varepsilon2\big(1+\ln(2/\varepsilon)\big)+O(\varepsilon^2).
\]
Substituting $\varepsilon=2e^{-2\eta}(1+O(e^{-2\eta}))$ gives $\varepsilon/2=e^{-2\eta}(1+O(e^{-2\eta}))$ and $\ln(2/\varepsilon)=2\eta+O(e^{-2\eta})$, so
\[
S(\eta) = e^{-2\eta}\big(1+O(e^{-2\eta})\big)\big(1+2\eta+O(e^{-2\eta})\big) = (1+2\eta)e^{-2\eta}\big(1+O(e^{-2\eta})\big),
\]
and taking logarithms gives $\ln S(\eta)=-2\eta+\ln(1+2\eta)+O(e^{-2\eta})$.
\end{proof}

\begin{figure}[htbp]
\centering
\includegraphics[width=0.62\linewidth]{fig4_entropy_tail.ps}
\caption{The linearised entropy tail of Corollary \ref{cor:entropy-asymp}(b): $\ln S(\eta)-\ln(1+2\eta)$ against $\eta=\QLC(\rho)$, compared with the reference line $-2\eta$. The two curves are visually indistinguishable beyond $\eta\approx2$.}
\end{figure}

\subsection{The Bloch-Sphere Radius and Distance-Based Measures of Mixedness}

For $d=2$, OMC admits a direct geometric meaning as the radial coordinate of the Bloch ball, which in turn makes several standard distance-based measures of mixedness exactly linear in OMC -- in sharp contrast to the entropy, which by Proposition \ref{prop:qubit} is an exact but strictly nonlinear function of OMC.

\begin{proposition}[Bloch-sphere radius identity]\label{prop:bloch}
Write a qubit density matrix as $\rho=\tfrac12(I+\vec r\cdot\vec\sigma)$, where $\vec\sigma=(\sigma_x,\sigma_y,\sigma_z)$ are the Pauli matrices and $r:=|\vec r|\in[0,1)$ is the Bloch-vector length. Then
\[
\QMC(\rho) = r, \qquad \QLC(\rho)=\operatorname{arctanh}(r).
\]
\end{proposition}
\begin{proof}
By the Pauli identity $(\vec a\cdot\vec\sigma)(\vec b\cdot\vec\sigma)=(\vec a\cdot\vec b)I+i(\vec a\times\vec b)\cdot\vec\sigma$ with $\vec a=\vec b=\vec r$, we have $(\vec r\cdot\vec\sigma)^2=r^2I$. Hence the traceless Hermitian operator $\vec r\cdot\vec\sigma$ has eigenvalues $\pm r$, so $\rho$ has eigenvalues $(1\pm r)/2$. Thus $\kappa(\rho)=(1+r)/(1-r)$, and
\[
\QMC(\rho) = \frac{\kappa(\rho)-1}{\kappa(\rho)+1} = \frac{\frac{1+r}{1-r}-1}{\frac{1+r}{1-r}+1} = \frac{(1+r)-(1-r)}{(1+r)+(1-r)} = r,
\]
and $\QLC(\rho)=\operatorname{arctanh}(\QMC(\rho))=\operatorname{arctanh}(r)$.
\end{proof}

\begin{remark}[Extremal cases: pure states and projections]
The Bloch-ball picture gives two extremal cases:
\begin{enumerate}
\item Pure states: $r\to1^-$ gives $\QMC(\rho)\to1$ and $\QLC(\rho)\to\infty$, consistent with the singular convention $\QLC=+\infty$ for singular positive operators.
\item Projections: A rank-1 projection $P=|v\rangle\langle v|$ has eigenvalues $\{1,0,0,\ldots\}$, hence $\kappa(P)=\infty$ and $\QMC(P)=1$, $\QLC(P)=+\infty$. More generally, any projection $P$ (rank $\le\dim H$) is singular, so $\QMC(P)=1$ and $\QLC(P)=+\infty$ by definition. Thus projections form the boundary of the positive cone where contrast is maximal.
\end{enumerate}
\end{remark}

\begin{corollary}[Linear geometric measures of mixedness]\label{cor:bloch-linear}
With $r=\QMC(\rho)$ as in Proposition \ref{prop:bloch},
\[
D(\rho,I/2) := \tfrac12\|\rho-I/2\|_1 = \frac r2, \qquad
\|\rho-I/2\|_2 = \frac r{\sqrt2}, \qquad
\operatorname{Tr}(\rho^2) = \frac{1+r^2}2.
\]
Thus the trace distance and the Hilbert--Schmidt distance to the maximally mixed state are exactly linear in $\QMC(\rho)$, while the purity is exactly linear in $\QMC(\rho)^2$.
\end{corollary}
\begin{proof}
$\rho-I/2=\tfrac12\vec r\cdot\vec\sigma$ has eigenvalues $\pm r/2$, so its trace norm (sum of absolute eigenvalues) is $r$, giving $D(\rho,I/2)=r/2$; its Hilbert--Schmidt norm is $\sqrt{(r/2)^2+(r/2)^2}=r/\sqrt2$. For the purity, using $(\vec r\cdot\vec\sigma)^2=r^2I$ from Proposition \ref{prop:bloch},
\[
\rho^2=\tfrac14\big(I+2\vec r\cdot\vec\sigma+(\vec r\cdot\vec\sigma)^2\big)=\tfrac14\big((1+r^2)I+2\vec r\cdot\vec\sigma\big),
\]
and taking the trace, using $\operatorname{Tr}(\vec r\cdot\vec\sigma)=0$ and $\operatorname{Tr}I=2$, gives $\operatorname{Tr}(\rho^2)=(1+r^2)/2$.
\end{proof}

\begin{remark}
Proposition \ref{prop:bloch} sharpens the picture of Remark \ref{rem:qubit-remark}: for $d=2$, OMC is not merely in exact bijection with $S$, it is literally the Euclidean radial coordinate of the Bloch ball, while $\QLC=\operatorname{arctanh}(r)$ is the corresponding hyperbolic (Thompson) radial coordinate furnished by Theorem \ref{thm:thompson} applied to the ray $[I]$ representing the maximally mixed state. This makes concrete, for the qubit case, the connection between the geometric identification of Section 3 and the hyperbolic geometry of the Bloch ball studied via Hilbert/Thompson-type projective metrics on qubit density matrices in \cite{park2015}. The linear behaviour of Corollary \ref{cor:bloch-linear} is special to these Euclidean/Hilbert--Schmidt-type distances and to $d=2$; it does not extend to the von Neumann entropy (Proposition \ref{prop:qubit} gives a nonlinear, though exactly bijective, relation), nor -- for the reasons of Section \ref{sec:nonbij} -- to any of these quantities once $d\ge3$. A natural follow-up, not pursued here, is whether the purity $\operatorname{Tr}(\rho^2)$ admits an analogue of the entropy-spread analysis of Section \ref{sec:nonbij} at fixed $\kappa$ for $d=3$; since purity is convex rather than concave along the same one-parameter family, the extremal structure would need to be re-derived rather than simply transplanted, and we leave this for future work (see also Section \ref{sec:open}).
\end{remark}

\subsection{Composite Systems: A Qutrit--Qubit Example}

Proposition \ref{prop:tensor-transfer} already shows that equality in subadditivity for pure tensor products of factors of dimension $\le3$ always forces commutativity, regardless of total ambient dimension. We record the qutrit--qubit instance explicitly.

\begin{example}[Qutrit $\otimes$ qubit density matrices]
Let $H=\mathbb{C}^3\otimes\mathbb{C}^2$, $d=6$. Let $\rho_2=\operatorname{diag}(2,1)/3\in M_2(\mathbb{C})_{++}$ and $\sigma_2\in M_2(\mathbb{C})_{++}$ be any positive matrix not diagonal in the same basis, chosen (as in the $n=2$ case of Theorem \ref{thm:threshold}) so that $\kappa(\rho_2\sigma_2)=\kappa(\rho_2)\kappa(\sigma_2)$. Fix any $\rho_1\in M_3(\mathbb{C})_{++}$ and set $\rho=\rho_1\otimes\rho_2$, $\sigma=\rho_1\otimes\sigma_2$ (rescaled to trace 1). By Proposition \ref{prop:tensor-transfer}, equality in subadditivity for $\rho,\sigma$ is governed entirely by the $M_2$ factor; since $n_2=2\le3$, any such equality forces $\rho_2\sigma_2=\sigma_2\rho_2$, hence $\rho\sigma=\sigma\rho$. No noncommuting equality-saturating pair exists at $d=6$ as long as both tensor legs stay $\le3$-dimensional, despite $d=6$ being comfortably above the $n=4$ threshold where Theorem \ref{thm:threshold} alone gives no such guarantee. This example is thus a direct corollary of Proposition \ref{prop:tensor-transfer} rather than an independent construction.
\end{example}

\subsection{Infinite-Dimensional Perspective: Spectral Measures and Weak Limits}\label{sec:inf-dim}

The entropy and purity spreads (Propositions \ref{prop:spread} and \ref{prop:purity-spread}) are stated for finite-dimensional density matrices. However, the underlying phenomenon generalizes to infinite dimensions via spectral measure theory.

For a positive trace-class operator $\rho\in B(H)_+$ with $\operatorname{Tr}(\rho)=1$ (density operator on infinite-dimensional $H$), write $\rho=\int_{\sigma(\rho)}\lambda\,dE_\lambda(\rho)$ where $dE_\lambda(\rho)$ is the spectral measure. The von Neumann entropy is $S(\rho)=-\operatorname{Tr}(\rho\log\rho)=\int_{\sigma(\rho)}s(\lambda)\,d\operatorname{Tr}(E_\lambda(\rho))$, where $s(\lambda)=-\lambda\ln\lambda$ is the pointwise entropy. If $\rho$ has discrete spectrum (eigenvalues $\{\lambda_i\}$ with multiplicities $m_i$), then
\[
S(\rho)=-\sum_im_i\lambda_i\ln\lambda_i.
\]

For a fixed condition number $\kappa(\rho)=\lambda_{\max}/\lambda_{\min}$, the set of density operators with that condition number forms a (possibly infinite-dimensional) convex set. Even though the full set may be infinite-dimensional, **the extremal spectra and entropy envelopes** that we computed for the finite-dimensional case extend to the closure of this set under weak-operator topology:

\begin{proposition}[Entropy envelope in infinite-dimensional limit]
Let $\{\rho_n\}_{n=1}^\infty$ be a sequence of finite-dimensional density matrices (on spaces of dimension $d_n\to\infty$) such that $\kappa(\rho_n)\to\kappa$ for some fixed $\kappa>1$. Then the cluster set of entropy values $\{S(\rho_n):n\ge1\}$ accumulates in the interval $[S_{\min}(\kappa),S_{\max}(\kappa)]$ (Proposition \ref{prop:spread}) as $d_n\to\infty$. Moreover, if $\rho_n\to\rho$ weakly in operator norm, where $\rho$ is a density operator on some Hilbert space (possibly infinite-dimensional) with $\kappa(\rho)=\kappa$, then $S(\rho)$ lies in the interval $[S_{\min}(\kappa),S_{\max}(\kappa))$.
\end{proposition}
\begin{proof}
The extremal spectra that saturate the upper and lower bounds for entropy (those achieving $S_P(\kappa)$ and $S_Q(\kappa)$ in Proposition \ref{prop:spread}) can be realized on arbitrarily large finite-dimensional spaces by padding with many copies of the extremal structure. Taking weak limits preserves the condition number and entropy (continuity of $S$ in the weak-operator topology on the set of trace-class operators with bounded spectrum), and the weak-limit density operator has entropy in the span of the cluster set.
\end{proof}

This observation shows that the finite-dimensional spreads are not artifacts of low dimension, but rather structural features that persist (and are actually extremal) in the infinite-dimensional limit.

\subsection{Non-bijectivity for $d\ge3$: The Qutrit Spread}\label{sec:nonbij}

For $d=3$, fix the condition number $\kappa=\lambda_1/\lambda_3>1$, where $\lambda_1\ge\lambda_2\ge\lambda_3>0$ are the eigenvalues of $\rho$. Writing $t=\lambda_1$, the trace constraint $\lambda_1+\lambda_2+\lambda_3=1$ together with $\lambda_3=t/\kappa$ forces
\[
\lambda_2(t) = 1-t\Big(1+\frac1\kappa\Big),
\]
and the ordering constraint $\lambda_3\le\lambda_2\le\lambda_1$ restricts $t$ to
\[
t \in \Big[\frac{\kappa}{2\kappa+1},\ \frac{\kappa}{\kappa+2}\Big],
\]
whose lower endpoint is the degeneracy $\lambda_1=\lambda_2$ and whose upper endpoint is $\lambda_2=\lambda_3$. This is a genuine one-parameter family of spectra sharing the same $\kappa$, unlike the qubit case, where $\kappa$ determines the spectrum outright.

\begin{proposition}[Entropy spread at fixed condition number, $d=3$]\label{prop:spread}
Along the family above, $S(t)=-\sum_i\lambda_i(t)\ln\lambda_i(t)$ is strictly concave in $t$. Consequently:
\begin{enumerate}[label=(\alph*)]
\item The maximum $S_{\max}(\kappa)$ is attained at the unique interior point where $\lambda_2/\lambda_1=\kappa^{-1/(\kappa+1)}$, equivalently at
\[
\lambda_1^* = \frac1{1+\kappa^{-1/(\kappa+1)}+1/\kappa}, \qquad \lambda_2^*=\kappa^{-1/(\kappa+1)}\lambda_1^*, \qquad \lambda_3^*=\lambda_1^*/\kappa.
\]
\item The minimum $S_{\min}(\kappa)$ is attained at one of the two boundary spectra,
\[
S_P(\kappa) = \frac{2\kappa}{2\kappa+1}\ln\frac{2\kappa+1}{\kappa} + \frac1{2\kappa+1}\ln(2\kappa+1) \qquad (\lambda_1=\lambda_2),
\]
\[
S_Q(\kappa) = \frac{\kappa}{\kappa+2}\ln\frac{\kappa+2}{\kappa} + \frac2{\kappa+2}\ln(\kappa+2) \qquad (\lambda_2=\lambda_3),
\]
and $S_Q(\kappa) < S_P(\kappa)$ for every $\kappa>1$ (proved analytically in Lemma \ref{lem:SPSQ} below), so $S_{\min}(\kappa)=S_Q(\kappa)$.
\end{enumerate}
Moreover $S_{\min}(\kappa)<S_{\max}(\kappa)$ strictly for every $\kappa>1$, so $\kappa$ (equivalently OMC or OLC) does not determine $S$ for $d=3$: the fiber over each value is a nondegenerate interval, not a point.
\end{proposition}
\begin{proof}
\emph{Concavity.} The map $t\mapsto(\lambda_1,\lambda_2,\lambda_3)(t)$ is affine, and Shannon entropy is strictly concave on the probability simplex; a strictly concave function composed with a non-constant affine map is strictly concave in $t$. A strictly concave function on a closed interval attains its maximum at an interior stationary point (if one exists) and its minimum only at an endpoint.

\emph{Stationary point.} Differentiating $S(t)=-t\ln t-\lambda_2\ln\lambda_2-\lambda_3\ln\lambda_3$ with $\lambda_2'(t)=-(\kappa+1)/\kappa$, $\lambda_3'(t)=1/\kappa$, and setting $dS/dt=0$, the constant terms cancel, leaving
\[
(\kappa+1)\ln\lambda_2 = \kappa\ln\lambda_1+\ln\lambda_3.
\]
Substituting $\lambda_3=\lambda_1/\kappa$ gives $\lambda_2^{\kappa+1}=\lambda_1^{\kappa+1}/\kappa$, i.e.\ $\lambda_2/\lambda_1=\kappa^{-1/(\kappa+1)}$, from which $\lambda_1^*,\lambda_2^*,\lambda_3^*$ follow via the trace constraint.

\emph{Endpoints.} Direct substitution of $t=\kappa/(2\kappa+1)$ and $t=\kappa/(\kappa+2)$ into $S(t)$ gives $S_P,S_Q$ above; part (b) is proved in Lemma \ref{lem:SPSQ}.
\end{proof}

We now supply the analytic proof of $S_Q(\kappa)<S_P(\kappa)$ that Proposition \ref{prop:spread}(b) uses; the original argument for this inequality was only checked numerically, which we regard as insufficient for a paper otherwise proved throughout by closed-form estimates.

\begin{lemma}\label{lem:SPSQ}
For every $\kappa>1$, $S_Q(\kappa) < S_P(\kappa)$.
\end{lemma}
\begin{proof}
First simplify both closed forms. Writing $S_P(\kappa)=\tfrac{2\kappa}{2\kappa+1}[\ln(2\kappa+1)-\ln\kappa]+\tfrac1{2\kappa+1}\ln(2\kappa+1)$ and collecting the $\ln(2\kappa+1)$ terms (their coefficients sum to $\tfrac{2\kappa+1}{2\kappa+1}=1$) gives
\[
S_P(\kappa) = \ln(2\kappa+1) - \frac{2\kappa}{2\kappa+1}\ln\kappa.
\]
An identical manipulation gives
\[
S_Q(\kappa) = \ln(\kappa+2) - \frac{\kappa}{\kappa+2}\ln\kappa.
\]
Define $\varphi(\kappa):=S_P(\kappa)-S_Q(\kappa)$ for $\kappa\ge1$. Then
\[
\varphi(\kappa) = \ln\frac{2\kappa+1}{\kappa+2} - \ln\kappa\left(\frac{2\kappa}{2\kappa+1}-\frac{\kappa}{\kappa+2}\right).
\]
A direct computation gives
\[
\frac{2\kappa}{2\kappa+1}-\frac{\kappa}{\kappa+2} = \kappa\cdot\frac{2(\kappa+2)-(2\kappa+1)}{(2\kappa+1)(\kappa+2)} = \frac{3\kappa}{(2\kappa+1)(\kappa+2)},
\]
so, writing $D(\kappa):=(2\kappa+1)(\kappa+2)=2\kappa^2+5\kappa+2$,
\[
\varphi(\kappa) = \ln\frac{2\kappa+1}{\kappa+2} - \frac{3\kappa\ln\kappa}{D(\kappa)}.
\]
At $\kappa=1$: $\tfrac{2\kappa+1}{\kappa+2}=1$ so the first term vanishes, and $\ln\kappa=0$ kills the second term; hence $\varphi(1)=0$ (consistent with $S_P(1)=S_Q(1)=\ln3$, both computed directly from the simplified forms above).

We show $\varphi'(\kappa)>0$ for $\kappa>1$, which together with $\varphi(1)=0$ gives $\varphi(\kappa)>0$, i.e.\ $S_P(\kappa)>S_Q(\kappa)$, for all $\kappa>1$. Differentiating the first term,
\[
\frac{d}{d\kappa}\ln\frac{2\kappa+1}{\kappa+2} = \frac2{2\kappa+1}-\frac1{\kappa+2} = \frac{2(\kappa+2)-(2\kappa+1)}{D(\kappa)} = \frac{3}{D(\kappa)}.
\]
Differentiating the second term, with $D'(\kappa)=4\kappa+5$,
\[
\frac{d}{d\kappa}\left(\frac{3\kappa\ln\kappa}{D(\kappa)}\right) = \frac{3(\ln\kappa+1)}{D(\kappa)} - \frac{3\kappa\ln\kappa\,D'(\kappa)}{D(\kappa)^2}.
\]
Hence
\[
\varphi'(\kappa) = \frac{3}{D(\kappa)} - \frac{3(\ln\kappa+1)}{D(\kappa)} + \frac{3\kappa\ln\kappa\,D'(\kappa)}{D(\kappa)^2}
= \frac{-3\ln\kappa}{D(\kappa)} + \frac{3\kappa(4\kappa+5)\ln\kappa}{D(\kappa)^2}
= \frac{3\ln\kappa\big[\kappa(4\kappa+5)-D(\kappa)\big]}{D(\kappa)^2}.
\]
Now $\kappa(4\kappa+5)-D(\kappa) = (4\kappa^2+5\kappa) - (2\kappa^2+5\kappa+2) = 2\kappa^2-2 = 2(\kappa-1)(\kappa+1)$, so
\[
\varphi'(\kappa) = \frac{6(\kappa-1)(\kappa+1)\ln\kappa}{D(\kappa)^2}.
\]
For $\kappa>1$, each of $(\kappa-1)$, $(\kappa+1)$, $\ln\kappa$, and $D(\kappa)^2$ is (strictly) positive, so $\varphi'(\kappa)>0$ on $(1,\infty)$. Since $\varphi$ is continuous on $[1,\infty)$ with $\varphi(1)=0$ and $\varphi'>0$ on $(1,\infty)$, we conclude $\varphi(\kappa)>0$, i.e.\ $S_Q(\kappa)<S_P(\kappa)$, for every $\kappa>1$.
\end{proof}

\begin{corollary}\label{cor:smax-bounds}
For every $\kappa>1$, $\ln2 \le S_{\max}(\kappa) < \ln3$.
\end{corollary}
\begin{proof}
As $\kappa\to\infty$, $\kappa^{-1/(\kappa+1)}\to1$, so $\lambda_1^*\to\tfrac1{1+1+0}=\tfrac12$, $\lambda_2^*\to\tfrac12$, $\lambda_3^*\to0$; hence $S_{\max}(\kappa)\to\ln2$ from above, the convergence being monotone decreasing in $\kappa$ (Figure \ref{fig:entropy-spread} and Table 1). As $\kappa\to1^+$, $S_{\max}(\kappa)\to\ln3$ (strictly, since $\kappa>1$). Combined with $S_{\max}(\kappa)>\ln2$ for every finite $\kappa$ (the limit $\ln2$ is never attained), this gives $\ln2\le S_{\max}(\kappa)<\ln3$ for all $\kappa>1$.
\end{proof}

\begin{remark}
As $\kappa\to1^+$ the interval collapses: $S_{\min},S_{\max}\to\ln3$. As $\kappa\to\infty$, $S_{\min}(\kappa)\to0$ while the entropy-maximising spectrum degenerates to $(\tfrac12,\tfrac12,0)$ -- the extremal $d=3$ configuration effectively reduces to a maximally mixed qubit with one eigenvalue vanishing, which is exactly the content of Corollary \ref{cor:smax-bounds}. The spread $S_{\max}(\kappa)-S_{\min}(\kappa)$ is monotonically increasing in $\kappa$, so the failure of bijectivity gets \emph{worse}, not better, as states move away from maximally mixed.
\end{remark}

\begin{figure}[htbp]
\centering
\includegraphics[width=0.72\linewidth]{fig3_entropy_spread.ps}
\caption{Entropy spread $[S_{\min}(\eta),S_{\max}(\eta)]$ for $d=3$ at fixed $\eta=\mathrm{OLC}(\rho)=\tfrac12\ln\kappa$ (Proposition \ref{prop:spread}), $\eta\in[0,4]$, linear scale. $S_{\max}(\eta)\to\ln2$ and $S_{\min}(\eta)\to0$ as $\eta\to\infty$, illustrating Corollary \ref{cor:smax-bounds}.}
\label{fig:entropy-spread}
\end{figure}

\begin{center}
\begin{tabular}{c|ccccccccc}
$\kappa$ & 1 & 4 & 9 & 19 & 50 & 100 & 500 & 1000 & 10000 \\
\hline
$S_{\max}(\kappa)$ & 1.0986 & 0.9745 & 0.8689 & 0.7974 & 0.7424 & 0.7213 & 0.7004 & 0.6972 & 0.6937 \\
$S_{\min}(\kappa)$ & 1.0986 & 0.8676 & 0.6002 & 0.3805 & 0.1897 & 0.1101 & 0.0288 & 0.0158 & 0.0020
\end{tabular}

\smallskip
Table 1: Selected values underlying Figure \ref{fig:entropy-spread}, cross-checked against the closed forms of Lemma \ref{lem:SPSQ}.
\end{center}

\subsection{Purity Spread for $d=3$}\label{sec:purity-spread}

Corollary \ref{cor:bloch-linear} already gives, for $d=2$, an exact closed-form linear relation between purity and $\QMC(\rho)^2$. For $d=3$ the same one-parameter family of spectra used in Section \ref{sec:nonbij} lets us carry out the analogous analysis for the purity $\operatorname{Tr}(\rho^2)=\lambda_1^2+\lambda_2^2+\lambda_3^2$ at fixed condition number $\kappa$. Unlike the entropy case, this reduces to an honest quadratic in $t$, so every extremal quantity below is a rational function of $\kappa$ -- there is no transcendental equation to solve -- and, since the parabola opens upward rather than downward, the roles of the interior and boundary extrema are exactly reversed relative to Proposition \ref{prop:spread}.

\begin{proposition}[Purity spread at fixed condition number, $d=3$]\label{prop:purity-spread}
Along the family $\lambda_1=t$, $\lambda_2(t)=1-t(1+1/\kappa)$, $\lambda_3=t/\kappa$, $t\in\big[\tfrac{\kappa}{2\kappa+1},\tfrac{\kappa}{\kappa+2}\big]$ of Section \ref{sec:nonbij}, the purity $P(t):=\lambda_1^2+\lambda_2^2+\lambda_3^2$ is a strictly convex quadratic function of $t$. Consequently:
\begin{enumerate}[label=(\alph*)]
\item The minimum is attained at the unique interior stationary point $t_*=\dfrac{\kappa(\kappa+1)}{2(\kappa^2+\kappa+1)}$, which lies strictly inside the interval for every $\kappa>1$, and equals
\[
P_{\min}(\kappa) = \frac{\kappa^2+1}{2(\kappa^2+\kappa+1)}.
\]
\item The maximum is attained at one of the two boundary spectra,
\[
P_P(\kappa) = \frac{2\kappa^2+1}{(2\kappa+1)^2} \quad(\lambda_1=\lambda_2), \qquad
P_Q(\kappa) = \frac{\kappa^2+2}{(\kappa+2)^2} \quad(\lambda_2=\lambda_3),
\]
and $P_Q(\kappa) > P_P(\kappa)$ for every $\kappa>1$, so $P_{\max}(\kappa)=P_Q(\kappa)$.
\end{enumerate}
Moreover $1/3 < P_{\min}(\kappa) < 1/2$ and $1/3 < P_{\max}(\kappa) < 1$ for every $\kappa>1$; at $\kappa=1$ the interval degenerates to $t=1/3$ and $P_{\min}=P_{\max}=1/3$, the purity of the maximally mixed state $I/3$.
\end{proposition}
\begin{proof}
Write $a:=1+1/\kappa$, so $P(t)=t^2+(1-at)^2+(t/\kappa)^2=At^2-2at+1$ with $A:=1+a^2+1/\kappa^2=2(\kappa^2+\kappa+1)/\kappa^2>0$; this is a strictly convex (upward) parabola, so it attains its minimum at the vertex $t_*=a/A=\kappa(\kappa+1)/(2(\kappa^2+\kappa+1))$ and its maximum at one of the two endpoints of any closed subinterval.

\emph{Interior location of $t_*$.} We check $t_*>\kappa/(2\kappa+1)$: cross-multiplying by the positive denominators and dividing by $\kappa>0$, this is equivalent to $(\kappa+1)(2\kappa+1)>2(\kappa^2+\kappa+1)$, i.e.\ $2\kappa^2+3\kappa+1>2\kappa^2+2\kappa+2$, i.e.\ $\kappa>1$. Likewise $t_*<\kappa/(\kappa+2)$ is equivalent, after the same manipulation, to $(\kappa+1)(\kappa+2)<2(\kappa^2+\kappa+1)$, i.e.\ $\kappa^2+3\kappa+2<2\kappa^2+2\kappa+2$, i.e.\ $\kappa(\kappa-1)>0$, again equivalent to $\kappa>1$. So $t_*$ lies strictly inside the interval exactly when $\kappa>1$.

\emph{Minimum value.} For an upward parabola $At^2-2at+1$, the vertex value is $1-a^2/A$. Since $a^2/A=(\kappa+1)^2/(2(\kappa^2+\kappa+1))$,
\[
P_{\min}(\kappa) = 1-\frac{(\kappa+1)^2}{2(\kappa^2+\kappa+1)} = \frac{2(\kappa^2+\kappa+1)-(\kappa+1)^2}{2(\kappa^2+\kappa+1)} = \frac{\kappa^2+1}{2(\kappa^2+\kappa+1)}.
\]

\emph{Boundary values.} At $t=\kappa/(2\kappa+1)$ ($\lambda_1=\lambda_2=t$, $\lambda_3=t/\kappa$), $P=2t^2+t^2/\kappa^2=t^2(2+1/\kappa^2)$, which simplifies to $P_P(\kappa)=(2\kappa^2+1)/(2\kappa+1)^2$. At $t=\kappa/(\kappa+2)$ ($\lambda_1=t$, $\lambda_2=\lambda_3=t/\kappa$), $P=t^2+2t^2/\kappa^2=t^2(1+2/\kappa^2)$, which simplifies to $P_Q(\kappa)=(\kappa^2+2)/(\kappa+2)^2$.

\emph{$P_Q>P_P$ for $\kappa>1$.} Clearing denominators,
\[
P_Q(\kappa)-P_P(\kappa) = \frac{(\kappa^2+2)(2\kappa+1)^2-(2\kappa^2+1)(\kappa+2)^2}{(\kappa+2)^2(2\kappa+1)^2}.
\]
The numerator expands to $2\kappa^4-4\kappa^3+4\kappa-2=2(\kappa^4-2\kappa^3+2\kappa-1)$, and $\kappa^4-2\kappa^3+2\kappa-1$ factors as $(\kappa-1)(\kappa^3-\kappa^2-\kappa+1)=(\kappa-1)(\kappa-1)^2(\kappa+1)=(\kappa-1)^3(\kappa+1)$. Hence
\[
P_Q(\kappa)-P_P(\kappa) = \frac{2(\kappa-1)^3(\kappa+1)}{(\kappa+2)^2(2\kappa+1)^2},
\]
which is strictly positive for $\kappa>1$ (and zero only at $\kappa=1$), proving $P_{\max}(\kappa)=P_Q(\kappa)$.

\emph{Bounds.} $P_Q(\kappa)<1$ iff $\kappa^2+2<(\kappa+2)^2$, i.e.\ $\kappa>-1/2$, always true; $P_Q(\kappa)>1/3$ iff $3(\kappa^2+2)>(\kappa+2)^2$, i.e.\ $2(\kappa-1)^2>0$, true for $\kappa\neq1$. Similarly $P_{\min}(\kappa)<1/2$ iff $\kappa^2+1<\kappa^2+\kappa+1$, always true for $\kappa>0$, and $P_{\min}(\kappa)>1/3$ iff $3(\kappa^2+1)>2(\kappa^2+\kappa+1)$, i.e.\ $(\kappa-1)^2>0$, true for $\kappa\neq1$.
\end{proof}

\begin{remark}
Proposition \ref{prop:purity-spread} mirrors Proposition \ref{prop:spread} exactly, but with every qualitative feature reversed: purity is convex rather than concave along the family, so the extremum at the interior stationary point is a minimum rather than a maximum, and the two boundary spectra swap roles ($P_{\max}=P_Q$ here, versus $S_{\min}=S_Q$ there). No transcendental equation of the form $\lambda_2/\lambda_1=\kappa^{-1/(\kappa+1)}$ appears; every quantity above is a rational function of $\kappa$, and the inequality $P_Q>P_P$ is established by exact polynomial factorisation rather than by a calculus argument, which is possible precisely because purity is quadratic where entropy is logarithmic.
\end{remark}

\subsection{Fourier-Theoretic Explanation of Non-Bijectivity}

The existence of the entropy spread (Proposition \ref{prop:spread}) and purity spread (Proposition \ref{prop:purity-spread}) for $d=3$ — and their absence for $d=2$ (Proposition \ref{prop:qubit}) — has a natural explanation in terms of log-spectral Fourier transforms. This observation, while not requiring new proofs, provides the structural reason why the bijectivity breaks at $d=3$.

For any invertible positive operator $A\in B(H)_{++}$, define the log-spectral measure $\nu_A$ as the normalized push-forward of the spectral measure of $\log A$ to the log-space coordinate $\lambda\mapsto\ln\lambda$. That is, if $A=\sum_i\lambda_i P_i$ with $\lambda_i>0$ and orthogonal projections $P_i$, then
\[
\nu_A := \frac1d\sum_i\delta(\cdot-\ln\lambda_i),
\]
where $d=\dim H$ (or formally, the measure is $(1/\operatorname{Tr}I)\sum_i\operatorname{Tr}(P_i)\delta(\cdot-\ln\lambda_i)$ in infinite dimensions). The Fourier transform of this measure is
\[
\Phi_A(t) := \int_{-\infty}^{\infty} e^{it\lambda}d\nu_A(\lambda) = \frac1d\sum_i e^{it\ln\lambda_i} = \frac1d\operatorname{Tr}(e^{it\log A}).
\]

\begin{proposition}[OLC as spectral diameter; Fourier characterization for qubits]\label{prop:fourier-char}
For any $A\in B(H)_{++}$,
\[
\QLC(A) = \tfrac12\operatorname{diam}(\sigma_{\log}(A)),
\]
where $\sigma_{\log}(A):=\{\ln\lambda:\lambda\in\sigma(A)\}$ is the log-spectrum. For $d=2$ (qubits), the log-spectrum is symmetric about zero with support $[-q,+q]$ where $q=\QLC(A)$, hence $\nu_A=\tfrac12[\delta(-q)+\delta(+q)]$ and
\[
\Phi_A(t) = \tfrac12[e^{itq}+e^{-itq}] = \cos(tq).
\]
Thus the Fourier transform $\Phi_A$ is completely determined by the single parameter $q=\QLC(A)$: different qubits with the same OLC have identical Fourier transforms of their log-spectra.
\end{proposition}
\begin{proof}
$\QLC(A)=\tfrac12\ln(\lambda_{\max}/\lambda_{\min})=\tfrac12[\ln\lambda_{\max}-\ln\lambda_{\min}]=\tfrac12\operatorname{diam}(\sigma_{\log}(A))$ by definition. For $d=2$, the eigenvalues are $(1+r)/2$ and $(1-r)/2$ where $r=\QMC(A)$, so $\ln\lambda_i=\ln\frac{1\pm r}2$, and $\tfrac12[\ln\frac{1+r}{2}-\ln\frac{1-r}{2}]=\tfrac12\ln\frac{1+r}{1-r}=\operatorname{arctanh}(r)=q=\QLC(A)$. The log-spectrum support is thus $[\ln\frac{1-r}{2},\ln\frac{1+r}{2}]$, which is symmetric about its midpoint $\ln(1/2)$ iff we work in shifted coordinates $\lambda'=\lambda/(1/2)$, i.e. $\ln\lambda'=\ln\lambda-\ln(1/2)$, giving support $[-q,+q]$. The Fourier formula follows immediately.
\end{proof}

\begin{corollary}[Why $d\ge3$ admits multiple spectra per OLC]\label{cor:fourier-nonbij}
For $d\ge3$, multiple spectra with the same condition number $\kappa$ (equivalently, the same OLC $q$) can yield \emph{different} log-spectral measures $\nu_A$, hence different Fourier transforms $\Phi_A(t)$. This is the fundamental obstruction to a bijection between OLC and entropy (or purity) for $d\ge3$:
\begin{itemize}
\item[\textit{d=2}:] Every density matrix with $\kappa(\rho)$ fixed is determined up to unitary by the pair $(\lambda_{\max},\lambda_{\min})$, which determines both $q=\QLC(\rho)$ and $S(\rho)$ uniquely. The Fourier measure $\Phi_\rho(t)=\cos(tq)$ is OLC-determined, and $S$ is a function of $\Phi$ (hence of OLC alone).
\item[\textit{d=3}:] For fixed $\kappa$, there is a one-parameter family of eigenvalue triples $(\lambda_1,\lambda_2,\lambda_3)$ satisfying the trace and ordering constraints (Section \ref{sec:nonbij}). These triples have the same $\kappa$ (hence the same $q$), but their log-spectra have different geometries. For instance, $(\lambda_1,\lambda_1,\lambda_3)$ with $\lambda_3=\lambda_1/\kappa$ gives $\nu_{\text{P}}$ supported on $\{\ln\lambda_1, \ln\lambda_1/\kappa\}$ (a pair of masses), while $(\lambda_1,\lambda_3,\lambda_3)$ with $\lambda_3=\lambda_1/\kappa$ gives $\nu_{\text{Q}}$ supported on $\{\ln\lambda_1, \ln\lambda_1/\kappa, \ln\lambda_1/\kappa\}$ (two masses at different points). These two measures have different Fourier transforms, hence define different density matrices $\rho_{\text{P}},\rho_{\text{Q}}$ with $\kappa(\rho_{\text{P}})=\kappa(\rho_{\text{Q}})$ but $S(\rho_{\text{P}})\neq S(\rho_{\text{Q}})$. The spread $[S_{\min}(\kappa),S_{\max}(\kappa)]$ (Proposition \ref{prop:spread}) is the set of all entropies achievable by spectra with condition number $\kappa$, and it is wider than a single point because the Fourier measure is not OLC-determined.
\end{itemize}
\end{corollary}

This Fourier perspective makes transparent why Propositions \ref{prop:spread} and \ref{prop:purity-spread} are not curiosities but structural consequences: the non-bijectivity of OLC and entropy/purity for $d\ge3$ is because the log-spectral measure — and hence its Fourier transform — contains information beyond OLC. For $d=2$, that information is trivial (the measure is determined by its diameter). For $d\ge3$, it is not.

\begin{figure}[htbp]
\centering
\includegraphics[width=0.95\linewidth]{fig7a_entropy_envelope.ps}
\caption{\textbf{Entropy envelope and spread width, vs OLC (Prop. 5.12).} Left: admissible range $[S_{\min}(\eta), S_{\max}(\eta)]$ for three-level systems, with bounds $\ln 2 \le S_{\max} < \ln 3$, plotted directly against $\eta=\mathrm{OLC}(\rho)=\tfrac12\ln\kappa$ rather than $\kappa$ itself. Right: width of the entropy spread $S_{\max}-S_{\min}$ as a function of $\eta$, showing growing ambiguity — for large $\eta$ (high purity), there is more structural freedom in how that purity is distributed among the three levels.}
\label{fig:entropy_envelope}
\end{figure}

\begin{figure}[htbp]
\centering
\includegraphics[width=0.95\linewidth]{fig7b_entropy_margins.ps}
\caption{\textbf{Entropy margin and ratio analysis, vs OLC (Prop. 5.12).} Left: margin of the envelope from the fixed value $\ln 2$, showing how far $S_{\max}$ lies above and $S_{\min}$ lies below this reference, against $\eta=\mathrm{OLC}(\rho)$. Both decay from small $\eta$ but have different rates. Right: ratio $S_{\max}/S_{\min}$ illustrating the multiplicative separation of upper and lower bounds, revealing the asymmetry in how entropy constraints operate in dimension 3.}
\label{fig:entropy_margins}
\end{figure}

\begin{figure}[htbp]
\centering
\includegraphics[width=0.95\linewidth]{fig7c_purity_envelope.ps}
\caption{\textbf{Purity envelope and spread width, vs OLC (Prop. 5.16).} Left: admissible range $[P_{\min}(\eta), P_{\max}(\eta)]$ for three-level systems with convex structure, plotted against $\eta=\mathrm{OLC}(\rho)$. Right: width of the purity spread $P_{\max}-P_{\min}$, which is monotone decreasing toward pure states. The opposite concavity structure compared to entropy (Fig. 7a) reflects the dual nature of purity versus entropy as state measures.}
\label{fig:purity_envelope}
\end{figure}

\begin{figure}[htbp]
\centering
\includegraphics[width=0.95\linewidth]{fig7d_purity_margins.ps}
\caption{\textbf{Purity margin and range evolution, vs OLC (Prop. 5.16).} Left: distance of the purity envelope from its natural boundaries ($P=1$ for pure states, $P=1/3$ for maximally mixed), both decaying exponentially with different rates as $\eta=\mathrm{OLC}(\rho)$ increases. Right: purity range width evolution showing the shrinking admissible region as $\eta$ grows. \emph{Key insight:} Entropy and purity spreads coexist precisely because multiple spectral configurations with identical $\kappa$ (and thus identical OLC) yield different entropy and purity values — the root of structural non-bijectivity for $d \ge 3$.}
\label{fig:purity_margins}
\end{figure}

\subsubsection{Why $d\ge3$ admits multiple spectra per OLC}

For $d\ge3$, multiple spectra with the same condition number $\kappa$ (equivalently, the same OLC $q$) can yield different density matrices with different entropies and purities. This is the fundamental obstruction to a bijection between OLC and entropy (or purity) for $d\ge3$. The explanation lies in the log-spectral Fourier transform $\Phi_A(t)=\tfrac1d\operatorname{Tr}(e^{it\log A})$ (Section 5.5):

For $d=2$ (qubits), every density matrix with condition number $\kappa$ is determined up to unitary by the pair $(r, \text{arbitrary phase})$, because the log-spectrum is forced to be symmetric $[-q, +q]$ with $q=\QLC(\rho)$. Hence $\Phi_A(t)=\cos(tq)$ is entirely OLC-determined, and entropy is a function of OLC alone (Proposition 5.4).

For $d\ge3$ (qutrits and higher), the spectrum can have many different internal geometric configurations while maintaining the same $\kappa$ and thus the same OLC. Different configurations yield different log-spectral measures and hence different Fourier transforms $\Phi_A(t)$, even though they have identical condition numbers. This multiplicity of spectral geometries is the root cause of the entropy and purity spreads (Propositions 5.12 and 5.16) for $d=3$.


\begin{corollary}[Near-pure tail of the purity spread]\label{cor:purity-tail}
As $\kappa\to\infty$ (equivalently $\eta:=\tfrac12\ln\kappa=\QLC\to\infty$),
\[
1-P_{\max}(\eta) = 4e^{-2\eta}\big(1+O(e^{-2\eta})\big), \qquad
\tfrac12-P_{\min}(\eta) = \tfrac12e^{-2\eta}\big(1+O(e^{-2\eta})\big),
\]
equivalently $\ln\big(1-P_{\max}(\eta)\big)=\ln4-2\eta+O(e^{-2\eta})$ and $\ln\big(\tfrac12-P_{\min}(\eta)\big)=-\ln2-2\eta+O(e^{-2\eta})$.
\end{corollary}
\begin{proof}
With $\kappa=e^{2\eta}$, $1-P_Q(\kappa)=\big[(\kappa+2)^2-(\kappa^2+2)\big]/(\kappa+2)^2=(4\kappa+2)/(\kappa+2)^2=(4/\kappa)\big(1+O(1/\kappa)\big)=4e^{-2\eta}\big(1+O(e^{-2\eta})\big)$. Similarly $\tfrac12-P_{\min}(\kappa)=\big[(\kappa^2+\kappa+1)-(\kappa^2+1)\big]/\big(2(\kappa^2+\kappa+1)\big)=\kappa/\big(2(\kappa^2+\kappa+1)\big)=\big(1/(2\kappa)\big)\big(1+O(1/\kappa)\big)=\tfrac12e^{-2\eta}\big(1+O(e^{-2\eta})\big)$. Taking logarithms gives the stated asymptotics.
\end{proof}

\begin{remark}\label{rem:purity-slope}
The slope $-2$ in Corollary \ref{cor:purity-tail} matches the slope $-2$ found for the entropy tail in Corollary \ref{cor:entropy-asymp}(b). This is not a coincidence: in both boundary families, the smallest eigenvalue $\lambda_3=t/\kappa$ scales as $\Theta(1/\kappa)=\Theta(e^{-2\eta})$ near the relevant boundary, and it is this common exponential rate -- not any special feature of entropy or purity individually -- that fixes the slope. We do not attempt to establish this as a general statement for arbitrary spectral functionals $\rho\mapsto-\operatorname{Tr}f(\rho)$ (e.g.\ R\'enyi entropies of general order); we record it as a natural question in Section \ref{sec:open}.
\end{remark}

\subsection{Entanglement Remarks}\label{rem:entanglement-section}

The noncommuting equality example from Theorem \ref{thm:threshold}, after normalisation to unit trace, gives two density matrices $\rho_X,\rho_Y$ with additive OLC but noncommuting. For $\rho_Y=Y/9$ with $Y$ as in the example, the partial transpose over the second qubit yields the matrix
\[
\rho_Y^{T_2} = \frac19\begin{pmatrix}4&0&0&1\\0&2&0&0\\0&0&2&0\\1&0&0&1\end{pmatrix}.
\]
Its eigenvalues are $\tfrac{5+\sqrt{13}}{18},\tfrac29,\tfrac29,\tfrac{5-\sqrt{13}}{18}$, all strictly positive. Hence the PPT criterion is satisfied; for two-qubit states this implies separability. Thus saturation of subadditivity does not require entanglement.

By Remark \ref{rem:tensor-irred} below, we can say more precisely where a genuinely noncommuting saturating example must live: it cannot be a pure tensor product of factors of dimension $\le3$, since Proposition \ref{prop:tensor-transfer} rules those out regardless of total dimension.

\begin{corollary}\label{rem:tensor-irred}
Any noncommuting equality-saturating pair of density matrices $\rho,\sigma$ -- entangled or not -- must be tensor-irreducible, in the sense that it cannot be written as $\rho=\rho_1\otimes\cdots\otimes\rho_k$, $\sigma=\sigma_1\otimes\cdots\otimes\sigma_k$ with every factor of dimension $\le3$; it must be built (as in Theorem \ref{thm:threshold}) from a block that does not itself decompose this way.
\end{corollary}
\begin{proof}
Immediate from Proposition \ref{prop:tensor-transfer}: any such decomposition with all $n_i\le3$ would force $\rho_i\sigma_i=\sigma_i\rho_i$ for every $i$, hence $\rho\sigma=\sigma\rho$.
\end{proof}

\begin{remark}\label{rem:entanglement}
As noted in Remark \ref{rem:polar}, it is tempting to try to extend Theorem \ref{thm:threshold} to general (not necessarily positive) invertible $X,Y$ by passing to absolute values $|X|,|Y|$, since $\kappa(X)=\kappa(|X|)$ and $\QLC(X)=\QLC(|X|)$. This is valid for statements purely about $\kappa$ or OLC, but not for the commutativity conclusion itself: $|X|$ and $|Y|$ commuting says nothing about whether $X$ and $Y$ commute (Pauli matrices $X=\sigma_x$, $Y=\sigma_y$ give $|X|=|Y|=I$ trivially commuting while $XY=-YX$). Theorem \ref{thm:threshold} genuinely requires $X,Y$ positive from the outset, since Lemma \ref{lem:align} uses the eigenspaces $\Emax,\Emin$ of $X,Y$ themselves. One can engineer entangled examples by choosing the common extremal eigenspaces to be entangled; we leave the explicit construction, constrained now by Proposition \ref{prop:tensor-transfer} and Corollary \ref{rem:tensor-irred}, for future work.
\end{remark}

\subsection{The Holevo Quantity and the Quantum Jensen--Shannon Divergence}\label{sec:qjsd}

For an ensemble $\{(p_i,\rho_i)\}_{i=1}^N$ of quantum states with average state $\bar\rho=\sum_i p_i\rho_i$, the Holevo $\chi$-quantity $\chi=S(\bar\rho)-\sum_i p_iS(\rho_i)$ upper-bounds the classically accessible information about the ensemble label under any measurement. For the equiprobable binary case $\{(\tfrac12,\rho_1),(\tfrac12,\rho_2)\}$, $\chi$ coincides with the quantum Jensen--Shannon divergence,
\[
\operatorname{QJSD}(\rho_1,\rho_2) := S(\bar\rho) - \tfrac12\big[S(\rho_1)+S(\rho_2)\big], \qquad \bar\rho=\tfrac12(\rho_1+\rho_2),
\]
a quantity already flagged, but not developed, as a keyword of our earlier paper \cite{abed2024}. We record here the exact qubit formula this connection yields, together with a genuinely dimension-free structural fact about $\bar\rho$; we also record, as an explicit warning to a future reader, why the most natural attempt to extend the qubit formula into a nontrivial higher-dimensional \emph{bound} on QJSD fails.

\begin{figure}[htbp]
\centering
\includegraphics[width=0.95\linewidth]{fig5b_qjsd_contours.ps}
\caption{\textbf{QJSD surface: 2D projections (Prop. 5.24).} Two complementary planar views of QJSD as a function of $\mathrm{OMC}(\rho_1)$ and the relative angle $\theta$, at fixed $\mathrm{OMC}(\rho_2)=0.6$. Left: contour plot with labeled levels, revealing non-uniform level spacing (denser near maximum where structure is richest). Right: heat map with alternative colormap (plasma) to emphasize the zone of high QJSD. Both panels show the same underlying landscape but highlight different structural features.}
\label{fig:qjsd_contours}
\end{figure}

\begin{figure}[htbp]
\centering
\includegraphics[width=0.95\linewidth]{fig5c_qjsd_slices.ps}
\caption{\textbf{QJSD cross-sections (Prop. 5.24).} Slices through the QJSD surface, at fixed $\mathrm{OMC}(\rho_2)=0.6$. Left: QJSD as a function of $\mathrm{OMC}(\rho_1)$ for five fixed angles ($\theta = 0, \pi/4, \pi/2, 3\pi/4, \pi$), showing monotone increase with $\mathrm{OMC}(\rho_1)$ and $\theta=\pi$ (orthogonal) always dominating. Right: QJSD as a function of $\theta$ for five fixed values of $\mathrm{OMC}(\rho_1)$ ($0.1, 0.3, 0.5, 0.7, 0.9$), showing dramatic rise toward $\theta=\pi$ (orthogonal states) and vanishing at $\theta=0$ (commuting states).}
\label{fig:qjsd_slices}
\end{figure}

\begin{figure}[htbp]
\centering
\includegraphics[width=0.95\linewidth]{fig5d_qjsd_ridges.ps}
\caption{\textbf{QJSD ridge analysis: max/min traces (Prop. 5.24), OLC view.} At fixed $\mathrm{OMC}(\rho_2)=0.6$. Left: along $\eta_1=\mathrm{OLC}(\rho_1)$ (the unbounded coordinate, as opposed to the bounded $\mathrm{OMC}(\rho_1)\in[0,1)$ used in Fig. 5c), the maximum QJSD (at $\theta=\pi$) grows toward $\ln 2$ as $\eta_1\to\infty$; the minimum (at $\theta=0$) grows more slowly, showing that orthogonal orientation dominates even at modest OLC. Right: along $\theta$, the maximum saturates at $\ln 2$ as $\theta \to \pi$ with $\mathrm{OMC}(\rho_1) \to 1$ (orthogonal pure states); minimum structure near $\theta=0$ scales with state mixture. Shaded regions show the admissible range at each parameter value.}
\label{fig:qjsd_ridges}
\end{figure}

\begin{proposition}[Exact qubit QJSD formula]\label{prop:qjsd-qubit}
Let $\rho_1,\rho_2$ be qubit density matrices with Bloch vectors $\vec r_1,\vec r_2$ of lengths $r_1,r_2\in[0,1)$ and angle $\theta$ between them, and write $H(x):=\ln2-\tfrac12\big[(1+x)\ln(1+x)+(1-x)\ln(1-x)\big]$ for the qubit entropy-vs-OMC function of Proposition \ref{prop:qubit} (extended by $H(1):=0$). Then
\[
\bar r := \QMC(\bar\rho) = \tfrac12\sqrt{r_1^2+r_2^2+2r_1r_2\cos\theta}, \qquad
\operatorname{QJSD}(\rho_1,\rho_2) = H(\bar r) - \tfrac12\big[H(r_1)+H(r_2)\big].
\]
\end{proposition}
\begin{proof}
The Bloch vector of $\bar\rho=\tfrac12(\rho_1+\rho_2)$ is $\tfrac12(\vec r_1+\vec r_2)$, so
\[
\bar r=\big|\tfrac12(\vec r_1+\vec r_2)\big|=\tfrac12\sqrt{|\vec r_1|^2+|\vec r_2|^2+2\vec r_1\cdot\vec r_2}=\tfrac12\sqrt{r_1^2+r_2^2+2r_1r_2\cos\theta}.
\]
By Proposition \ref{prop:qubit}, $S(\rho)=H(\QMC(\rho))$ for any qubit density matrix, and by Proposition \ref{prop:bloch}, $\QMC(\rho_1)=r_1$, $\QMC(\rho_2)=r_2$, $\QMC(\bar\rho)=\bar r$; substituting into the definition of QJSD gives the stated formula.
\end{proof}

\begin{remark}[Consistency checks]
Two extreme cases recover known values. If $\theta=0$ ($\rho_1,\rho_2$ commute, i.e.\ diagonal in a common basis), $\bar r=(r_1+r_2)/2$ and the formula reduces to the classical Jensen--Shannon divergence between the two Bernoulli distributions $(1\pm r_i)/2$, as it must for commuting states. If $\theta=\pi$ and $r_1=r_2=1$ (orthogonal pure states, e.g.\ $|0\rangle\langle0|$ and $|1\rangle\langle1|$), then $\bar r=0$ and $\operatorname{QJSD}=H(0)-\tfrac12[H(1)+H(1)]=\ln2-0=\ln2$, the known maximal value of QJSD for perfectly distinguishable states.
\end{remark}

\begin{remark}[Relation to the classical Michelson--JSD equivalence, and why it does not lift past $d=2$]\label{rem:bruni}
Bruni, Rossi, and Vitulano \cite{bruni2012} prove a formal equivalence, for classical two-level (binary) distributions, between the Michelson contrast of an image distortion and the Jensen--Shannon divergence it induces. Proposition \ref{prop:qjsd-qubit} is exactly the operator lift of that equivalence to $d=2$: $r_i=\QMC(\rho_i)$ plays the role of the Michelson contrast, and QJSD is recovered as a closed-form function of $(r_1,r_2,\theta)$ alone. For $d=2$ this is a genuine, checkable correspondence, not a cosmetic analogy.

It does not extend to $d\ge3$, and the reason is exactly Proposition \ref{prop:spread}: for $d\ge3$, OLC (hence OMC) no longer determines the spectrum, so it cannot determine QJSD either. Concretely, fix $\kappa_1=5$ and take the two extremal $d=3$ spectra of Proposition \ref{prop:spread} sharing that condition number, $\lambda^{\min\text{-entropy}}=(5,1,1)/7$ and $\lambda^{\max\text{-entropy}}\approx(0.596,0.331,0.072)$ (both diagonal in the same basis, so no angle ambiguity enters); pair each against a fixed third state $\rho_2$ with spectrum $(3,1.7,1)/5.7$. Both choices of $\rho_1$ have identical $\QMC(\rho_1)$, identical $\QMC(\rho_2)$, and identical (trivial, commuting) relative orientation with $\rho_2$ -- yet direct computation gives $\operatorname{QJSD}=0.0220$ for the min-entropy spectrum against $0.0080$ for the max-entropy spectrum, a $\sim\!170\%$ relative difference. So already the most favorable case for a closed-form generalization -- commuting states, where there is no orientation angle to worry about -- rules one out: no function of $(\QMC(\rho_1),\QMC(\rho_2),\theta)$ alone can equal QJSD once $d\ge3$. (A structural reason sits underneath the numerics too: a $d$-level Bloch-type representation carries $d^2-2$ independent relative-orientation parameters between two states, not one, so $\theta$ alone is the wrong-shaped object to parametrize QJSD once $d\ge3$, quite apart from whether OMC pins down the spectrum.) The Bruni-type equivalence is therefore a $d=2$ phenomenon, not a defect of QJSD or of OMC individually -- it is a direct corollary of the same non-bijectivity already established in Proposition \ref{prop:spread}.
\end{remark}

\begin{figure}[htbp]
\centering
\includegraphics[width=0.95\linewidth]{fig6a_qjsd_contours.ps}
\caption{\textbf{QJSD contours: commuting vs orthogonal (Prop. 5.26).} Contour plots of QJSD over $\eta_1,\eta_2\in[0,1.5]$, the range where both cases actually vary (a shared $[0,\ln 2]$ scale would make the left panel look flat, since it never exceeds $0.14$ while the right panel reaches $0.50$ -- so each panel uses its own scale, noted in the axis label). Left: commuting case ($\theta=0$), true maximum $\approx0.14$ over this window. Right: orthogonal case ($\theta=\pi$), true maximum $\approx0.50$, approaching $\ln 2$ only as both $\eta_i\to\infty$.}
\label{fig:qjsd-contours}
\end{figure}

\begin{figure}[htbp]
\centering
\includegraphics[width=0.95\linewidth]{fig6b_qjsd_slices.ps}
\caption{\textbf{QJSD diagonal and cross-sections (Prop. 5.26).} Left: diagonal analysis ($\eta_1=\eta_2$) confirms QJSD($\rho,\rho$) = 0 by construction, and shows the maximum possible divergence when comparing a state to an orthogonal one. Right: cross-sections with fixed $\eta_1$ values (0.3, 0.9, 1.8) show QJSD increases monotonically with $\eta_2$, demonstrating how purity of the second state increases distinguishability.}
\label{fig:qjsd-slices}
\end{figure}

\begin{figure}[htbp]
\centering
\includegraphics[width=0.95\linewidth]{fig6c_qjsd_marginals.ps}
\caption{\textbf{QJSD marginal analysis and orthogonal gain (Prop. 5.26).} Left: range of QJSD over $\eta_2$ for each fixed $\eta_1$, showing growing dispersion as state 1 becomes purer — purer states are more sensitive to the properties of state 2. Right: difference map quantifying the gain in QJSD from orthogonal ($\theta=\pi$) versus commuting ($\theta=0$) orientation as a function of both condition numbers. \emph{Key insight:} QJSD is not determined by OLC values alone; it requires knowledge of both condition numbers and their relative orientation.}
\label{fig:qjsd-marginals}
\end{figure}



\begin{proposition}[Ensemble-averaging never worsens conditioning]\label{prop:ensemble-contraction}
Let $\{(p_i,\rho_i)\}_{i=1}^N$ be a finite ensemble of positive invertible operators $\rho_i\in B(H)_{++}$ with weights $p_i>0$, and $\bar\rho:=\sum_{i=1}^N p_i\rho_i$. Then
\[
\QLC(\bar\rho) \le \max_{1\le i\le N} \QLC(\rho_i).
\]
\end{proposition}
\begin{proof}
Apply the sum inequality (9) to $A_i:=p_i\rho_i>0$: $\QLC\big(\sum_iA_i\big)\le\max_i\QLC(A_i)$. Since OLC is invariant under positive scaling (Section 2.2), $\QLC(A_i)=\QLC(p_i\rho_i)=\QLC(\rho_i)$, and $\sum_iA_i=\bar\rho$.
\end{proof}

\begin{remark}[A bound that does not work, recorded to save a future attempt]\label{rem:qjsd-vacuous}
It is natural to try to combine Proposition \ref{prop:ensemble-contraction} with the entropy bound (15) to obtain a higher-dimensional bound on QJSD: for the binary ensemble, (15) gives $S(\bar\rho)\le\ln d+2\QLC(\bar\rho)$, and Proposition \ref{prop:ensemble-contraction} gives $\QLC(\bar\rho)\le\max\{\QLC(\rho_1),\QLC(\rho_2)\}$, so
\[
\operatorname{QJSD}(\rho_1,\rho_2) \le S(\bar\rho) \le \ln d + 2\max\{\QLC(\rho_1),\QLC(\rho_2)\}.
\]
Since $\QLC\ge0$ always, the right-hand side is at least $\ln d\ge\ln2$ for $d\ge2$, so this bound is never better than the already-known dimension-free bound $\operatorname{QJSD}\le\ln2$ -- it is vacuous for every ensemble, in every dimension. We record this explicitly so the natural-looking combination is not rediscovered as if it were new content; the genuine content of this section is the exact formula of Proposition \ref{prop:qjsd-qubit} for $d=2$ and the dimension-free structural fact of Proposition \ref{prop:ensemble-contraction}, not a nontrivial bound on QJSD for $d\ge3$. Figure \ref{fig:vacuous-gap} shows this vacuity by plotting the gap between the naive bound and the true bound: it is strictly positive for all $\eta\ge0$ and all $d\ge2$.
\end{remark}

\begin{figure}[htbp]
\centering
\includegraphics[width=0.95\linewidth]{fig10a_vacuous_gap.ps}
\caption{\textbf{Vacuity of the QJSD bound (Remark 5.29).} We define the gap $g(d,\eta):=[\ln d+2\eta]-\ln 2$ between the naive bound $\ln d+2\eta$ (obtained by combining Proposition \ref{prop:ensemble-contraction} with the entropy bound (15)) and the true, dimension-free bound $\ln 2$: $g>0$ means the naive bound is strictly worse than what was already known. Left: $g$ for each dimension $d=2,\ldots,9$ over $\eta\in[0,1.5]$; every curve stays strictly positive, so the naive bound is never competitive on this range (and $g$ only grows for larger $\eta$). Right: the same gap over the $(d,\eta)$ plane, linear color scale, restricted to $\eta\in[0,1.5]$ -- this is the window where $g$ is small enough for the gradient to be informative; for $\eta\gtrsim2$, $g$ already exceeds $3$ for every $d$ shown, and a wider window would just look uniformly saturated.}
\label{fig:vacuous-gap}
\end{figure}

\section{Structural Interpretations, Outlook, and Conclusion}

\subsection{Rapidity and Hyperbolic-Geometric Interpretation}

The relation $\QMC(A)=\tanh(\QLC(A))$ also reveals a structural connection with the hyperbolic parametrisation underlying the notion of rapidity in special relativity. In one spatial dimension, a future-directed unit timelike vector in Minkowski space can be written as $u(\eta)=(\cosh\eta,\sinh\eta)$, where $\eta$ is the rapidity, and $\beta=v/c=\tanh\eta$. For collinear boosts, rapidities are additive: $\eta_{12}=\eta_1+\eta_2$.

A formally analogous bounded-to-unbounded parametrisation occurs for OMC: $\QMC(A)=\tanh(\QLC(A))$, $\QLC(A)=\operatorname{arctanh}(\QMC(A))$. Thus OLC can be viewed as a rapidity-type coordinate for the bounded contrast variable OMC. It is important, however, to distinguish this rapidity analogy from the operator-geometric interpretation established earlier. For positive invertible operators, $\QLC(A)=\delta_T([I],[A])$, so OLC is exactly the projective Thompson distance from the identity ray -- a genuine metric identification on the positive cone. By contrast, the connection with rapidity arises from the scalar hyperbolic parametrisation $x=\tanh t$, $t=\operatorname{arctanh} x$, and from the analogy between the corresponding composition structures.

Consequently, the Thompson geometry and the rapidity geometry are distinct geometric structures. The former is associated with the order and spectral structure of the positive operator cone, whereas the latter is associated with Lorentzian geometry and hyperbolic angles in Minkowski space. The common appearance of $\tanh$ does not identify these geometries. This distinction is particularly relevant for the composition law: rapidity is additive for collinear Lorentz boosts, $\eta_{12}=\eta_1+\eta_2$, whereas OLC satisfies only the subadditivity estimate $\QLC(AB)\le\QLC(A)+\QLC(B)$. Exact rapidity additivity is a consequence of the one-dimensional Lorentz boost group, whereas OLC subadditivity follows from the submultiplicativity of the operator condition number; noncommutativity and possible misalignment of extremal spectral directions prevent a general additive law \cite{rindler2006}.

\subsection{Structural Remarks: Thompson Geometry and Harmonic Analysis}

The geometric identification of Theorem \ref{thm:thompson} (OLC as Thompson distance) and the Fourier-theoretic explanation of Corollary \ref{cor:fourier-nonbij} (Section 5.5) together situate OLC within two complementary analytical frameworks. The Thompson metric governs the contraction and ordering properties of OLC on the positive cone, while the log-spectral Fourier perspective (Section 5.5) explains why functionals of the spectrum (entropy, purity, divergences) are constrained but not determined by OLC in dimensions $\ge3$. These perspectives may deepen through future work connecting them via harmonic analysis on convex cones of log-spectral measures, but the structural foundations are already visible in the present results.

\subsection{Open Questions}\label{sec:open}
\begin{itemize}
\item Whether the common slope $-2$ found for both the entropy tail (Corollary \ref{cor:entropy-asymp}(b)) and the purity tail (Corollary \ref{cor:purity-tail}) near a pure qubit state extends to R\'enyi entropies $S_\alpha$ of every order $\alpha$, as the mechanism of Remark \ref{rem:purity-slope} suggests.
\item Connection between the tensor-irreducibility structure of Theorem \ref{thm:threshold} and the representation theory of SU(2) and SU(3): whether the $n=4$ threshold is linked to the finite-dimensionality of irreducible representations, and whether analogous thresholds exist for higher groups or other operator systems.
\item Explicit construction of entangled density matrices with additive OLC, constrained by Proposition \ref{prop:tensor-transfer} to be tensor-irreducible (Corollary \ref{rem:tensor-irred}).
\item Relationship between the equality condition of Theorem \ref{thm:threshold} and Knabe-type local-to-global gap criteria for 1D spin chains, given the structural analogy between the $n\le3$ vs.\ $n\ge4$ threshold and block-size criteria.
\item Rigorous treatment of Theorem \ref{thm:threshold} for the intermediate infinite-dimensional regime (compact perturbations of the identity, or operators whose spectrum accumulates only away from the extremal points), as sketched in Remark \ref{rem:infinite-dim}.
\end{itemize}

\subsection{Speculative Directions}

\subsubsection{Optical Systems and Channel Structure}

The dimensional threshold $n\le3$ vs.\ $n\ge4$ has an intriguing parallel in optical systems. The classical Michelson contrast was originally defined for optical interferometry with single-wavelength light. In modern imaging and sensing, several levels of channel multiplexing are routinely used:

\begin{enumerate}
\item **Color imaging (d=3)**: RGB channels are standard. The operator Michelson contrast applies directly to the color-intensity matrix $M_\mathbf{I}:=\operatorname{diag}(R,G,B)$, as noted in Remark 5.2. The $d=3$ threshold of Theorem \ref{thm:threshold} thus sits exactly at the dimensionality of human color vision.

\item **Multispectral and hyperspectral imaging (d >> 3)**: Modern satellite and aerial imaging systems often acquire $d=8$–$d=224$ spectral bands \cite{goetz1985}. For these high-dimensional systems, our finite-dimensional threshold predicts that noncommuting equality cases in subadditivity should exist generically, reflecting more complex spectral-structure phenomena than in the $d=3$ regime. The threshold between ``typical'' commutativity ($d\le3$) and ``anomalous'' noncommutativity ($d\ge4$) may thus have operational significance: for RGB systems, aligned extremal spectra are forced; for hyperspectral systems, they are not.

\item **Polarization and quantum optics (d=4, 6)**: Fully-general polarization states of light lie in a 4-dimensional space (Stokes parameters, or equivalently the coherency matrix). Higher-order optical systems (orbital angular momentum, spatial modes) can involve $d\ge4$ degrees of freedom. Multimode quantum optics with entangled photons naturally involves tensor products of mode spaces, making Proposition \ref{prop:tensor-transfer} (tensor-irreducibility of the threshold) directly applicable: even a 2-photon system with correlated modes in $\mathbb{C}^2\otimes\mathbb{C}^2$ (d=4 total) will exhibit the full $d=4$ noncommuting behavior if the entanglement structure is not a pure product of two $d=2$ subsystems \cite{raussendorf2003,o2003}.
\end{enumerate}

The observation that our mathematical threshold coincides with the dimensionality of practical optical channel counts (3 for color, 4+ for advanced systems) invites the question: does the commutativity-forcing property of $d\le3$ have a physical origin in the geometry of optical state spaces? This remains open. At minimum, the coincidence suggests that the dimensional threshold is not a mathematical curiosity but potentially reflects real structure in signal processing and optics.

\subsubsection{Representation Theory and Phase Structure}

The recurrence of the dimensions $2$ and $3$ as the exact threshold for forced commutativity (Theorem \ref{thm:threshold}) is intriguing, given their role as the dimensions of the fundamental representations of $SU(2)$ and $SU(3)$. However, we have identified no representation-theoretic mechanism rigorously linking this numerology to Theorem \ref{thm:threshold} or Proposition \ref{prop:tensor-transfer}. A key obstacle is that the \textbf{condition number is blind to phases}: for any unitary $U$, $\kappa(U)=1$ (spectrum on the unit circle), hence $\MC(U)=0$ and $\QLC(U)=0$. This means OLC and the subadditivity analysis capture only magnitude, not phase information. 

To probe the $n\le3$ vs.\ $n\ge4$ dimensional threshold through representation-theoretic properties, one would need to extend the analysis beyond condition numbers to the **phase structure of the spectrum**. This naturally suggests a Fourier-analytic approach: the log-spectral Fourier transform $\Phi_A(t)=\frac1d\operatorname{Tr}(e^{it\log A})$ (Section 5.5) encodes both magnitude (through the spectral radius $\sup|\Phi_A(t)|$) and phase (through oscillations and winding as $t$ varies). A representation-theoretic interpretation of the dimensional threshold would likely emerge from analyzing the phase structure of $\Phi_A(t)$ for basis representations of $SU(2), SU(3),\ldots$ and understanding how this phase structure changes at $n=4$. This remains an open direction.

\subsubsection{Computational Analogy}

The identity $\QMC(A)=\tanh(\QLC(A))$ admits a computational reading: with $z=\QLC(A)$, $\QMC(A)=\tanh z$ is formally the output of the hyperbolic tangent applied to an unbounded pre-activation coordinate, in the sense familiar from neural-network activation functions \cite{goodfellow2016}. This should be regarded purely as a naming analogy, not as a proposed architecture or as evidence of any deeper connection.

\subsection{Conclusion}
We have studied the logarithmic parametrisation of the operator Michelson contrast, obtaining the subadditive functional $\QLC(A)=\tfrac12\ln\kappa(A)$, which we call the Operator Logarithmic Contrast. This logarithmic lift is the operator analogue of the classical logarithmic contrast and fits into a commutative diagram with the Michelson contrast. We proved its subadditivity, geometric interpretation as a projective Thompson distance, an accompanying Lipschitz-continuity property that also pins down the correct limiting behaviour at singular operators, and analysed the equality case, showing that $n=4$ is the smallest dimension in which noncommuting equality cases can occur -- and that this threshold is fundamentally about tensor-irreducibility rather than raw ambient dimension: equality forces commutativity for any total dimension, provided the operators factor as pure tensor products of blocks of size at most $3$, via a fully spelled-out inductive argument. Additional properties include contraction, sum inequalities, trace estimates, and estimates for double-sided products $X^{1/2}YX^{1/2}$. Applied to quantum information, these results yield an exact closed-form bijection between OMC and von Neumann entropy for qubits, purity bounds at fixed condition number for qutrits, and a proof -- established by complete analytic arguments rather than numerical checks -- that no bijection between OMC and entropy exists for $d\ge3$, together with explicit envelopes on the resulting ambiguity. The logarithmic parametrisation reveals a rich structure inherent in the operator Michelson contrast: not only as a geometric quantity (Thompson distance, Theorem \ref{thm:thompson}), but also as a spectral bandwidth determined by log-spectral distributions (Proposition \ref{prop:fourier-char} and Corollary \ref{cor:fourier-nonbij}, Section 5.5). These foundational perspectives invite future development of finer Fourier-theoretic properties, connections to commutator bounds and Baker--Campbell--Hausdorff theory, and applications to related operator-theoretic problems where the interplay between condition numbers and spectral structure plays a central role.


\begin{thebibliography}{99}

\bibitem{abed2024} S. A. Abed, I. Nikolaeva, A. Novikov, ``Generalisation of Michelson contrast for operators and its properties'', \emph{Lobachevskii Journal of Mathematics} 45(8) (2024), 3835--3848, DOI 10.1134/S1995080224603400.

\bibitem{bruni2012} V. Bruni, E. Rossi, D. Vitulano, ``On the Equivalence Between Jensen--Shannon Divergence and Michelson Contrast'', \emph{IEEE Trans. Inform. Theory} 58(7) (2012), 4278--4288, DOI 10.1109/TIT.2012.2191472.

\bibitem{jourlin1988} M. Jourlin, J. C. Pinoli, ``A model for logarithmic image processing'', \emph{Journal of Microscopy} 149(1) (1988), 21--35.

\bibitem{jourlin1989} M. Jourlin, J. C. Pinoli, R. Zeboudj, ``Contrast definition and contour detection for logarithmic images'', \emph{Journal of Microscopy} 156(1) (1989), 33--40.

\bibitem{dixmier1969} J. Dixmier, \emph{Les $C^*$-alg\`ebres et leurs repr\'esentations}, Gauthier-Villars, Paris (1969).

\bibitem{thompson1963} A. C. Thompson, ``On certain contraction mappings in a partially ordered vector space'', \emph{Proc. Amer. Math. Soc.} 14 (1963), 438--443, DOI 10.1090/S0002-9939-1963-0149237-7.

\bibitem{corach1994} G. Corach, H. Porta, L. Recht, ``Convexity of the geodesic distance on spaces of positive operators'', \emph{Illinois J. Math.} 38(1) (1994), 87--94, DOI 10.1215/ijm/1255986889.

\bibitem{lawson2001} J. D. Lawson, Y. Lim, ``The geometric mean, matrices, metrics, and more'', \emph{Amer. Math. Monthly} 108(9) (2001), 797--812.

\bibitem{beneduci2014} R. Beneduci, L. Moln\'ar, ``On the standard K-loop structure of positive invertible elements in a $C^*$-algebra'', \emph{J. Math. Anal. Appl.} 420 (2014), 551--562, DOI 10.1016/j.jmaa.2014.05.009.

\bibitem{abe2015} T. Abe, O. Hatori, ``Generalized gyrovector spaces and a Mazur--Ulam theorem'', \emph{Publ. Math. Debrecen} 87 (2015), 393--413, DOI 10.5486/PMD.2015.7234.

\bibitem{park2015} J. Park, S. Kim, ``Hilbert Projective Metric on a Gyrogroup of Qubit Density Matrices'', \emph{Rep. Math. Phys.} 76(3) (2015), 389--400, DOI 10.1016/S0034-4877(15)30039-2.

\bibitem{kim2016} S. Kim, ``Gyrovector Spaces on the Open Convex Cone of Positive Definite Matrices'', \emph{Math. Interdiscip. Res.} 1(1) (2016), 173--185.

\bibitem{rindler2006} W. Rindler, \emph{Relativity: Special, General, and Cosmological}, 2nd ed., Oxford University Press, Oxford (2006).

\bibitem{goodfellow2016} I. Goodfellow, Y. Bengio, A. Courville, \emph{Deep Learning}, MIT Press, Cambridge, MA (2016).

\bibitem{goetz1985} A. F. H. Goetz, G. Vane, J. E. Conel, ``Imaging spectrometry for Earth remote sensing'', \emph{Science} 228 (1985), 1147--1153, DOI 10.1126/science.228.4704.1147.

\bibitem{raussendorf2003} R. Raussendorf, H. J. Briegel, ``A one-way quantum computer'', \emph{Phys. Rev. Lett.} 86 (2001), 5188--5191, DOI 10.1103/PhysRevLett.86.5188.

\bibitem{o2003} M. J. O'Brien, ``Photonic quantum technologies'', \emph{Nature Photonics} 3 (2009), 687--695, DOI 10.1038/nphoton.2009.229.

\end{thebibliography}


\appendix
\section{Reproducing the Figures}\label{app:code}

Every figure in this paper is generated by the single self-contained script below. It has no dependencies beyond \texttt{numpy} and \texttt{matplotlib}, and writes each PNG to the current working directory under the filename referenced by the corresponding \verb|\includegraphics| command in the source, so it can be run as-is to regenerate the figure set. The script is organised into six parts: shared helper functions (binary entropy, the extremal $d=3$ spectra of Propositions \ref{prop:spread} and \ref{prop:purity-spread}, and the qubit QJSD formula of Proposition \ref{prop:qjsd-qubit}); Figures 1--4 (the qubit entropy--contrast relation and the $d=3$ entropy spread); Figures 5b--5d (the QJSD surface as a function of Bloch radius and angle); Figures 6a--6c (QJSD as a function of the two condition numbers $\eta_1,\eta_2$); Figures 7a--7d (entropy and purity spreads for $d=3$); and Figure 10a (the vacuity of the naive QJSD bound, Remark \ref{rem:qjsd-vacuous}). A final function reproduces the numerical counterexample of Remark \ref{rem:bruni} in isolation, printing the two QJSD values and their relative difference.

\textbf{Note:} The Python script that generated the figures is not included in the arXiv submission due to its length, but it is available from the authors upon request.

% The actual \lstinputlisting command is omitted to avoid compilation errors.
% \lstinputlisting[caption={\texttt{generate\_all\_figures.py}: reproduces every figure in the paper.}]{generate_all_figures.py}

\end{document}
